\documentclass[11pt,a4paper]{article}
\usepackage[margin=27mm]{geometry}
\usepackage[T1]{fontenc}
\usepackage{lmodern,amsmath,amssymb,amsthm,mathtools}
\usepackage{microtype,booktabs,longtable,array}
\usepackage[hyphens]{url}
\usepackage[colorlinks=true,linkcolor=blue,urlcolor=blue,citecolor=blue]{hyperref}
\newtheorem{theorem}{Theorem}[section]
\newtheorem{proposition}[theorem]{Proposition}
\newtheorem{lemma}[theorem]{Lemma}
\newtheorem{corollary}[theorem]{Corollary}
\theoremstyle{definition}
\newtheorem{definition}[theorem]{Definition}
\newtheorem{remark}[theorem]{Remark}
\newtheorem{example}[theorem]{Example}
\newcommand{\Z}{\mathbb Z}
\newcommand{\Q}{\mathbb Q}
\newcommand{\F}{\mathbb F}
\DeclareMathOperator{\ord}{ord}
\title{Exact selectors, character energy, and integral return\\
for the Erd\H{o}s--Straus equation}
\author{The Clankers}
\date{12 September 2026}
\begin{document}
\maketitle
\begin{abstract}
The bounded divisor formulation of the Erd\H{o}s--Straus equation admits
exact arithmetic, character-theoretic and geometric descriptions. We prove
completeness of the first-half denominator selector and give its inverse
coordinate maps. Inversion parity yields an exact collision-energy minimum
for a divisor distribution avoiding all solution targets; resolving this
minimum by quadratic characters strengthens the resulting existence test.
The original arithmetic gate is also the exact order of a specified affine
cocycle. An integral three-coordinate change identifies its cyclic action
with a hexagonal quotient, retaining the third coordinate and index-three
gluing. The homogeneous cyclotomic quartic is the same quadratic form on
positive squares, with a complete description of its dihedral square returns.
Local compactification maps are given using actual branch-cycle lengths.
These results do not establish universal Erd\H{o}s--Straus existence.
\end{abstract}
\tableofcontents
\clearpage
\section{Research lineage and conventions}
The Boolean, positivity, cyclotomic and prime-character directions in the
joint investigation are credited to JT as identified in the supplied
dialogue. The maintainer's programme asks for exact maps linking those
arithmetic conditions to the geometric structures, without discarding
coordinates or scales. Assistant-derived arguments are recorded as such,
not retroactively attributed to a human who requested their investigation.
The source locators and hashes in \texttt{../PROVENANCE.md} distinguish the
complete dialogue, recovered source packages, and the newly written proofs.

The type-I/type-II arithmetic has its literature antecedent in Elsholtz
and Tao \cite{et}; the received papers preserve their more detailed
source-specific references. The integral monodromy used below is taken
from the displayed matrices of the circulated period manuscript
\cite{period}, through the received torsion-cover note. Only those explicit
matrices enter these algebraic proofs. The global sphere claim of that
manuscript is not a premise. Permutation triples and branched covers have
the classical context described by Sijsling and Voight \cite{sv}.

Each section declares its own coordinates: in particular a primitive
tuple's scale $D$ and an affine system's level $D$ are named locally and
are not silently identified. Ordered denominators, channel marks,
divisor exponents and branch origins are retained. Finite computations
are distinguished from the general written proofs.

% Copy-ready mathematical module. Requires amsmath, amssymb, amsthm,
% hyperref and theorem/lemma/proposition/corollary environments.
% No repository files are modified by this intake artifact.
\section{The complete first-half selector and its cyclotomic relations}
\label{sec:bp-increment}

The source dialogue is JT and ChatGPT, \emph{Boolean product formulation},
9--11 September 2026, exported 12 September 2026
(\href{https://chatgpt.com/share/6aa445c7-f990-83ea-a2b9-0a9cae13dc52}{shared source dialogue}).
The preserved transcript has SHA-256
\texttt{19d54ae24d712598}\allowbreak
\texttt{c14f222ee81eb19f1}\allowbreak
\texttt{92c9e944fbbfe500}\allowbreak
\texttt{9f91780709792f3}.
The locators below refer to the one-based lines of that export.
JT supplies the Boolean-product, cyclotomic, generalized-congruence and
adversarial directions; the following arguments state the mathematical
claims with their exact scopes. They make no blanket claim of novelty.
In particular no universal Erd\H{o}s--Straus theorem or counterexample is
asserted. The first-half proof and the explicit complete permutation fibres
below fill out the relation between the retained full-range code and the
dialogue's first-half code.

Write \(\operatorname{ES}(p)\) for the existence of positive integers
\(x,y,z\) such that \(4/p=1/x+1/y+1/z\).
Prime factorization, quadratic reciprocity with its supplementary laws,
and cyclicity of the multiplicative group of a finite field are used in
their usual classical forms. Every additional arithmetic implication in
this section is proved here.

\subsection{The denominator bound, divisor maps and exact code}

\begin{theorem}[First-half denominator bound]\label{thm:bp-first-half}
Let \(p\) be prime and \(p\equiv1\pmod4\). The minimum denominator in
every positive integer Erd\H{o}s--Straus solution lies strictly between
\(p/4\) and \(p/2\).
\end{theorem}
\begin{proof}
Retain a permutation of the original coordinates and name its ordered
values \(x\le y\le z\). Positivity gives \(1/x<4/p\), so \(x>p/4\).
Suppose \(x>p/2\); equality is impossible because \(p\) is odd.
Then
\[
 \frac4p\le\frac1x+\frac2y<\frac2p+\frac2y,
\]
so \(y<p\). Neither \(x\) nor \(y\) is divisible by \(p\).
The integer equation \(4xyz=p(xy+xz+yz)\) gives \(p\mid z\).
Write \(z=pk\). The original equality becomes
\[
 (4k-1)xy=pk(x+y).
\]
Thus \(t=(4k-1)/p\) is a positive integer,
\(t\equiv3\pmod4\), \(\gcd(t,k)=1\), and \(txy=k(x+y)\).
Write \(x=g\alpha,y=g\beta\), where \(g=\gcd(x,y)\) and
\(\gcd(\alpha,\beta)=1\). Then
\[
 tg\alpha\beta=k(\alpha+\beta),\qquad
 \gcd(\alpha\beta,\alpha+\beta)=1.
\]
Hence \(k=\alpha\beta D\) for an integer \(D>0\), and
\[
 tg=D(\alpha+\beta),\quad \gcd(t,D)=1,\quad
 pt=4\alpha\beta D-1.
\]
The inequality \(x>p/2\) implies
\[
 2D\alpha(\alpha+\beta)>4\alpha\beta D-1,
 \qquad 2D\alpha(\alpha-\beta)>-1.
\]
The last left side is an integer, so it is nonnegative and
\(\alpha\ge\beta\). Applying the same argument to \(y>p/2\)
gives \(\beta\ge\alpha\). Coprimality forces
\(\alpha=\beta=1\). Now \(tg=2D\) and \(\gcd(t,D)=1\)
give \(t\mid2\), contrary to \(t\equiv3\pmod4\).
The retained permutation recovers the original coordinate order.
\end{proof}

Fix henceforth a prime \(p=12h+1\), \(h\ge1\), and put
\[
 a=3h+j+1,\quad R=4a-p=4j+3,\quad S=pa,
 \qquad 0\le j<3h.
\]
Thus \(p/4<a<p/2\), \(3\le R\le p-2\), and
\(\gcd(R,pa)=1\). Indeed \(a<p\) is prime to \(p\), a common
divisor of \(a,R\) divides \(p\), and a common divisor of \(p,R\)
divides \(4a\); also \(R\) is odd. Define
\[
 E_a=\{u>0:u\mid a^2,\ R\mid4u+1\},\qquad
 M_a=\{u>0:u\mid a^2,\ R\mid4u+p\}.
\]
The middle condition is equivalently \(R\mid u+a\).

\begin{theorem}[Original ordered solutions and the bounded divisor code]
\label{thm:bp-code}
For \(u\mid a^2\), set
\[
 g=\gcd(a,u),\quad A=a/g,\quad B=u/g,\quad D=g^2/u.
\]
These are positive integers and define a bijection
\[
 \operatorname{Div}(a^2)\longleftrightarrow
 \{(A,B,D)>0:a=ABD,\ \gcd(A,B)=1\},\qquad u=B^2D.
\]
For \(\varepsilon\in\{0,1\}\), its hit condition is exactly
\[
 R\mid4u+p^\varepsilon
 \quad\Longleftrightarrow\quad R\mid A+p^{1-\varepsilon}B.
\]
On that set, define
\[
 C=\frac{A+p^{1-\varepsilon}B}{R},\qquad
 (x,y,z)=(ABD,p^\varepsilon ACD,pBCD).
 \tag{BP1}
\]
Then \(C>0\) is integral, \(4ABCD=A+p^{1-\varepsilon}B+pC\),
and (BP1) is an original ordered solution.

Put
\[
 M_H=6h,\quad J_H=3h,\quad\Lambda_H=36h^2,
 \quad N_H=216h^3=(p-1)^3/8.
\]
The map
\[
 c_H=(A-1)\Lambda_H+(p-1)j+2(B-1)+\varepsilon
 \tag{BP2}
\]
is a bijection from the rectangle
\(1\le A,B\le M_H,\ 0\le j<J_H,\ \varepsilon\in\{0,1\}\)
to \(0\le c_H<N_H\). On its decoded coordinates set
\[
 \chi_{H,p}(c_H)=
 \left\lfloor\frac{\gcd(4AB,p+4j+3)}{4AB}\right\rfloor
 \left\lfloor\frac1{\gcd(A,B)}\right\rfloor
 \left\lfloor\frac{\gcd(4j+3,A+p^{1-\varepsilon}B)}{4j+3}\right\rfloor.
\]
Then
\[
 \operatorname{ES}(p) \Longleftrightarrow\
 \exists\,0\le c_H<N_H:\chi_{H,p}(c_H)=1.
 \tag{BP3}
\]
\end{theorem}
\begin{proof}
Write \(a=gA,u=gB\) with \(\gcd(A,B)=1\). The divisibility
\(u\mid a^2\) implies \(B\mid gA^2\), hence \(B\mid g\).
Therefore \(D=g/B\) is integral and \(a=ABD,u=B^2D\).
Conversely these equalities give \(u\mid a^2\) with quotient
\(A^2D\), and \(\gcd(a,u)=BD\). This proves both inverse
compositions. Also \(A,B\le a\le M_H\).
Using \(4ABD=p+R\), one has
\[
 A(4u+p^\varepsilon)
 \equiv pB+p^\varepsilon A
 =p^\varepsilon(A+p^{1-\varepsilon}B)\pmod R.
\]
Both \(A\) and \(p\) are units modulo \(R\), proving the
biconditional. On its passing locus, multiplying the definition of
\(C\) by \(R=4ABD-p\) gives the tuple equation. The sum of the
three reciprocals in (BP1), over denominator \(pABCD\), has numerator
\(pC+p^{1-\varepsilon}B+A=4ABCD\).

For the code inverse divide \(c_H\) by \(\Lambda_H\) to recover
\(A-1\) and remainder \(v\), divide \(v\) by \(p-1\) to recover
\(j\) and remainder \(r\), then divide \(r\) by 2 to recover
\(B-1\) and \(\varepsilon\). The ranges of successive remainders
are precisely the prescribed ranges. The gcd-floor factors equal
the divisibility or coprimality bits because their quotients are
positive and at most one. Thus a passing code gives (BP1).

For the converse, Theorem~\ref{thm:bp-first-half} selects a first
denominator \(a\) in this interval, retaining its original position.
For the other denominators define
\[
 d_y=Ry-S=Sy/z>0,\quad d_z=Rz-S>0,
 \qquad d_y d_z=S^2.
\]
Conversely a positive \(d_y\mid S^2\) with \(R\mid d_y+S\)
determines \(d_z=S^2/d_y\) and
\(y=(d_y+S)/R,z=(d_z+S)/R\); the second quotient is integral
because \(d_y\equiv-S\) is a unit modulo \(R\).
Since \(p\nmid a\), the exponent \(i=v_p(d_y)\) is 0,1 or 2.
For \(i=0\), set \(\varepsilon=0,u=a^2/d_y\); multiplying
\(d_y+pa\equiv0\) by \(4u/a\) gives \(p(1+4u)\equiv0\).
For \(i=2\), set \(\varepsilon=0,u=d_y/p^2\); multiplication
by \(4/p\) gives the same congruence, with the final two original
denominators exchanged in (BP1). For \(i=1\), write \(d_y=pv\)
and set \(\varepsilon=1,u=a^2/v\). Then \(v\equiv-a\), hence
\(u\equiv-a\), giving the middle condition. All three cases
therefore give a passing code. The exterior exchange is recorded,
not identified with the original ordering. This proves (BP3).
\end{proof}

\begin{proposition}[Radix embedding and the full permutation fibres]
\label{prop:bp-radix}
Retain the repository's full rectangle
\[
 J_F=6h,\quad M_F=9h,\quad\Lambda_F=108h^2,
 \quad N_F=972h^3,
\]
with the same shell formula and code
\(c_F=(A-1)\Lambda_F+18hj+2(B-1)+\varepsilon\).
The full-range hits satisfying \(a\le6h\) are bijective to the
first-half hits by decoding and recoding the same four coordinates.
The map preserves \(a,R,S,A,B,C,D,u\), the ordered triple (BP1),
and every predicate on these coordinates, and
\[
 c_F-c_H=2(A-1)\Lambda_H+6hj.
 \tag{BP4}
\]
All other full-range hits still have a computable image among
first-half hits after retaining the coordinate permutation. More
precisely, let \(\mathcal H_{\min}\) consist of first-half hit codes
whose first denominator is minimal and whose middle-channel lift has
\(y<z\). There is a surjection from full-range hits to
\(\mathcal H_{\min}\). For a fixed code in \(\mathcal H_{\min}\),
its entire inverse fibre is obtained by the following finite rule:
take all distinct permutations \((b,c,d)\) of its lifted triple,
retain those with
\[
 3h+1\le b\le9h,\qquad
 v_p((4b-p)c-pb)\in\{0,1\},
 \tag{BP5}
\]
and recover their full codes by the inverse in the proof.
Thus the two existential selectors agree, although their full hit
sets, shell labels and total counts need not agree.
\end{proposition}
\begin{proof}
For a full-range hit with \(a\le6h\), the conditions \(AB\mid a\)
give \(A,B\le6h\), and \(j<3h\). The hit tests on those same
coordinates are identical. The two proved radix decoders are inverse,
and the values \(\Lambda_F=3\Lambda_H\), \(18h-12h=6h\)
give (BP4).

For the second statement lift a full code to its ordered triple.
Choose its minimum denominator as first coordinate. Theorem
\ref{thm:bp-first-half} puts it in the first-half interval. The
residual-divisor proof in Theorem~\ref{thm:bp-code} gives a unique
canonical exterior orientation, with residual exponent 0 for the
second denominator, or a middle orientation with both exponents 1.
In the latter case choose the smaller of the final two denominators
as second. They cannot be equal: equality would give
\(d_y=d_z=S\), hence \(R\mid2S\), and the unit condition would
give \(R\mid2\), impossible. This yields a unique code in
\(\mathcal H_{\min}\). In the exterior case exactly one of the
two residual exponents is 0 and the other 2, so its orientation is
also unique. The minimum value and the multiset of remaining values
are unique, even if a value occurs more than once; hence these rules
give a unique resulting code.

For a retained permutation in (BP5), put \(R_b=4b-p,S_b=pb\) and
\(d_c=R_bc-S_b\). When \(v_p(d_c)=0\), recover
\(\varepsilon=0,u=b^2/d_c\); when \(v_p(d_c)=1\), recover
\(\varepsilon=1,u=pb^2/d_c\). The proof of Theorem
\ref{thm:bp-code} gives the positive integers \(A,B,C,D\) and
the code with \(j=b-3h-1\). Since \(A,B\le b\le9h\), it is
in the full rectangle and its ordered lift is exactly \((b,c,d)\).
Conversely every full code in the inverse fibre has that triple as
one of the distinct permutations, its first coordinate is in the
full interval, and its canonical residual exponent is 0 or 1.
Thus (BP5) gives every point of the fibre and no duplicates.
Surjectivity follows because each minimal first-half code has its
identical full-radix lift.
\end{proof}

\begin{corollary}[Boolean products and the exact two-layer inclusion]
Define
\[
 S(p)=1-\prod_{c=0}^{N_H-1}(1-\chi_{H,p}(c)),\qquad
 S_2(p)=1-\prod_{c=0}^{2\Lambda_H-1}(1-\chi_{H,p}(c)).
\]
Then \(S_2(p)\le S(p)\), \(S(p)=1\) if and only if
\(\operatorname{ES}(p)\), and \(S_2(p)=1\) if and only if a
first-half hit has \(A\in\{1,2\}\). For any prime set \(\mathcal P\)
define its infinite Boolean product as the infimum of its finite
initial products. Then
\[
 \prod_{p\in\mathcal P}S(p)=1
 \quad\Longleftrightarrow\quad
 (\forall p\in\mathcal P)\operatorname{ES}(p).
\]
\end{corollary}
\begin{proof}
Each complement factor is a bit. Its finite product is one exactly
when every hit bit vanishes. The radix intervals for fixed \(A\)
are \([(A-1)\Lambda_H,A\Lambda_H-1]\), so the first two are
exactly the displayed smaller range. Inclusion gives the inequality.
The finite initial products form a nonincreasing sequence of bits;
one zero factor makes every later product zero, and absent such a
factor every finite product is one. Apply (BP3).
\end{proof}

These statements correspond to transcript lines 14--163, 180--387,
5203--5293 and 6465--6606. The repository's previous exact radix map
is Theorem \texttt{raw-first-half-code} in
\nolinkurl{raw_scale_hit_incidence.tex}; the broader full-range
bound and code occur in \nolinkurl{bounded_transport.tex} and
\nolinkurl{unary_boolean_transport.tex}.

\subsection{Endpoint forcing and the corrected generalized carrier}

\begin{theorem}[The endpoint factor theorem]\label{thm:bp-endpoint}
Let \(p\equiv1\pmod{24}\) be prime. An \(A=1\) hit with
\(R^2\le p\) and \(B=1\) or \(D=1\) exists if and only if
\((p+1)/2\) or \(p+4\) has a prime divisor congruent to 3 modulo 4.
If \(q\) is the least such divisor of the first integer, the
exterior code is \((p-1)(q-3)/4\). If it is the least such divisor
of the second integer, the middle code is \((p+1)(q-1)/4\).
\end{theorem}
\begin{proof}
Both integers are odd and 1 modulo 4. If either has a prime factor
3 modulo 4, the number of such prime factors counted with multiplicity
is positive and even. The least one \(q\) therefore has square at
most that integer. For \((p+1)/2\), this gives \(q^2<p\).
For \(p+4\), it gives \(q^2\le p+4\); since \(q^2-p\)
is divisible by 8, it follows that \(q^2\le p\).

In the first case use
\[
 (R,A,B,D,\varepsilon)=
 (q,1,1,(p+q)/4,0),\quad C=(p+1)/q.
\]
In the second case use
\[
 (R,A,B,D,\varepsilon)=
 (q,1,(p+q)/4,1,1),\quad C=(p+q+4)/(4q).
\]
All these quantities are positive integers; in the second case oddness
of \(q\) and \(q\mid p+4\) give the last integrality. Substitution
in (BP2) gives the stated codes, and the tuple equations prove the hits.

Conversely \(A=B=1\) makes the middle channel impossible because
\(R\nmid2\); the exterior channel requires \(R\mid p+1\).
If \(A=D=1\), then \(B=a\), and the middle channel requires
\(R\mid p+4\). The exterior channel would require
\(R\mid p^2+4\). A positive integer \(R\equiv3\pmod4\)
has a prime factor \(q\equiv3\pmod4\). Such a prime cannot divide
\(p^2+4\): division by the unit 4 would make \(-1\) a square
modulo \(q\). A square root of \(-1\) would have order 4 in a
group of order \(q-1\equiv2\pmod4\), impossible. In each remaining
case that prime \(q\mid R\) divides \((p+1)/2\) or \(p+4\).
\end{proof}

The factor hypothesis is not universal: at \(p=193\), the two
integers are the primes 97 and 197, both 1 modulo 4. Nevertheless
\((A,B,C,D,\varepsilon)=(1,5,138,10,0)\) gives the code 200 and
the solution \((50,1380,1331700)\). This retains the actual scope of
transcript lines 2570--2944 and its correction at 5149--5189.

\begin{theorem}[Exact generalized carrier]
\label{thm:bp-crt}
For \(u>0\) set
\(\kappa(u)=\prod_{\ell^e\parallel u}\ell^{\lceil e/2\rceil}\).
Let \(n>0\) and \(R>0\), \(R\equiv3\pmod4\). The system
\[
 t\equiv1\pmod{12n},\quad
 t\equiv-R\pmod{4\kappa(u)},\quad
 t\equiv-R-4u\pmod{4R}
 \tag{BP6}
\]
has an integer solution if and only if
\[
 \gcd(3n,\kappa(u))\mid(R+1)/4,\qquad
 \gcd(3n,R)\mid(R+4u+1)/4.
 \tag{BP7}
\]
Its complete solution set, when nonempty, is one class modulo
\(4\operatorname{lcm}(3n,\kappa(u),R)\).
For the exterior carrier retain \(R\mid4u+1\) and omit the third
congruence; only the first divisibility in (BP7) remains.
For a prescribed prime \(p\), actual membership in the carrier,
\(R<p\), and the original common divisor must be retained.
\end{theorem}
\begin{proof}
Prime valuations give \(u\mid a^2\) if and only if
\(\kappa(u)\mid a\), and \(\kappa(u)\mid u\). With
\(4a=p+R\), these are exactly
\[
 u\mid a^2\ \Longleftrightarrow\ p\equiv-R\pmod{4\kappa(u)},
 \quad
 R\mid a+u\ \Longleftrightarrow\ p\equiv-R-4u\pmod{4R}.
\]
For a finite system of integer congruences, pairwise residue agreement
modulo the corresponding gcds is necessary by subtraction. To prove
sufficiency, for each prime select a modulus having maximal exponent
of that prime. Its prescribed residue agrees with every other at all
lower powers by the pairwise condition. Ordinary CRT over these distinct
prime powers gives a solution. Differences of solutions are exactly
multiples of the lcm. Applying this criterion to (BP6), the first two
pair conditions are (BP7); the third is
\(\gcd(\kappa(u),R)\mid u\), automatic because \(\kappa(u)\mid u\).
The lcm gives the stated period. The same argument with two moduli
gives the exterior statement.

For a class \(b+m\mathbb Z\), intersection with an integer interval
\([\alpha,\beta]\) occurs exactly when
\[
 \left\lceil\frac{\alpha-b}{m}\right\rceil
 \le\left\lfloor\frac{\beta-b}{m}\right\rfloor.
\]
At \(\alpha=\beta=p\), this is precisely membership of \(p\).
Finally a single \(u\mid a^2\) has the unique exponent vector
\((e_\ell)\in\prod_{\ell\mid a}\{0,\ldots,2v_\ell(a)\}\).
Reduction modulo different prime powers dividing \(R\) uses this same
vector. Its complete fibre is the intersection of the local fibres
on that box, not their Cartesian product.
\end{proof}

At \(n=2\), (BP7) is exactly
\[
 2\mid u\Rightarrow R\equiv7\pmod8,\quad
 3\mid u\Rightarrow R\equiv2\pmod3,\quad
 3\mid R\Rightarrow u\equiv2\pmod3.
\]
Thus a modulo-8 condition cannot be removed: \(R=3,u=2\) gives
incompatible requirements \(t\equiv1\pmod{24}\) and
\(t\equiv5\pmod8\). Also \(R=7,u=1\) gives the compatible
class \(73\pmod{168}\), which does not contain 193. These are the
corrections of transcript lines 4740--5058 and the generalization at
6104--6158.

\subsection{Cyclotomic order, infinitude, coverage and a smaller cofactor}

\begin{theorem}[Cyclotomic Euclid and complete-family coverage]
\label{thm:bp-cyclotomic}
Let \(m=12n\), \(n>0\), and \(P>0\). Then
\[
 \Phi_m(mP)>1,\qquad \gcd(mP,\Phi_m(mP))=1,
\]
and every prime divisor of this value is 1 modulo \(m\).
There are infinitely many primes in that class. For each such prime
\(p\), there are exactly \(\varphi(m)/2\) integers
\(1\le r\le(p-1)/2\) with \(p\mid\Phi_m(mr)\).
Consequently, with \(P_T=\prod_{r=1}^T\Phi_m(mr)\),
\[
 p\equiv1\pmod m,\quad p\le2T+1\quad\Longrightarrow\quad p\mid P_T.
 \tag{BP8}
\]
In the Boolean convention of the preceding corollary,
\[
 \prod_{p\equiv1\ (m)}S(p)
 =\inf_{T\ge1}\prod_{r=1}^T
   \prod_{\substack{q\mid\Phi_m(mr)\\q\ \text{prime}}}S(q).
 \tag{BP9}
\]
This identity does not assert that either side equals one.
\end{theorem}
\begin{proof}
Use \(X^m-1=\prod_{d\mid m}\Phi_d(X)\) and \(\Phi_m(0)=1\).
For a prime \(q\nmid mr\), a root \(r\) of \(\Phi_m\)
has order exactly \(m\): a smaller order \(e\mid m\) would make
it a common root with a factor \(\Phi_d\), \(d\mid e<m\),
and hence a repeated root of \(X^m-1\). This contradicts the
nonzero derivative \(mr^{m-1}\) modulo \(q\). Conversely an
element of order \(m\) vanishes on none of the proper-divisor
factors, so it vanishes on \(\Phi_m\).

For \(X=mP\), the constant term gives \(\Phi_m(X)\equiv1\pmod X\),
so each prime divisor is prime to \(mX\). Its order is \(m\),
which divides \(q-1\). Pair primitive roots into nonreal conjugates.
For real \(X\ge2\), every pair contributes \(|X-\zeta|^2\)
strictly between \((X-1)^2\) and \((X+1)^2\); hence
\[
 (X-1)^{\varphi(m)}<\Phi_m(X)<(X+1)^{\varphi(m)}.
\]
This proves \(\Phi_m(mP)>1\). If \(P\) is the product of any
finite list of primes 1 modulo \(m\), a prime divisor of that value
is outside the list and in the same class, proving infinitude.

For a prescribed prime \(p\equiv1\pmod m\), the cyclic group
\(\mathbb F_p^\times\) has exactly \(\varphi(m)\) elements
of order \(m\). Multiplication by the unit \(m^{-1}\) gives exactly
that many roots of \(\Phi_m(mX)\). As \(4\mid m\), negation
preserves primitive \(m\)-th roots: the exponent of a primitive
root changes from a unit \(k\pmod m\) to \(k+m/2\), still a
unit at every odd prime and still odd. Thus \(\Phi_m\) is even.
The nonzero representatives pair as \(r,p-r\), giving the exact
half-interval count and in fact
\[
 p^{\varphi(m)/2}\mid\prod_{r=1}^{(p-1)/2}\Phi_m(mr).
\]
This proves (BP8). Every prime occurring in the right side of (BP9)
is in the required class, and (BP8) eventually includes every member
of that class. Repetitions have no effect because \(S(q)^e=S(q)\)
for \(e\ge1\). Both infima therefore vanish for exactly the same
zero factors and otherwise equal one.
\end{proof}

The polynomial change needed for the larger modulus is exact:
\[
 \Phi_{12}(X^n)
 =\prod_{\substack{d\mid n\\\gcd(n/d,6)=1}}\Phi_{12d}(X),
 \quad
 \Phi_{12n}(X)=
 \prod_{\substack{d\mid n\\\gcd(d,6)=1}}
 (X^{4n/d}-X^{2n/d}+1)^{\mu(d)}.
 \tag{BP10}
\]
Indeed \(\operatorname{ord}(\zeta^n)=\operatorname{ord}(\zeta)/
\gcd(\operatorname{ord}(\zeta),n)=12\) precisely when
\(\operatorname{ord}(\zeta)=12d\), \(d\mid n\), and
\(\gcd(n/d,6)=1\). Simple roots give the first polynomial
identity; multiplicative M\"obius inversion gives the second.
Thus \(\Phi_{12n}(X)=X^{4n}-X^{2n}+1\) exactly when
\(n=2^\alpha3^\beta\). Any other prime divisor of \(n\)
gives an additional nonconstant factor in the first identity.
Changing only the argument does not impose order \(12n\):
\(\Phi_{12}(24)=13\cdot73\cdot349\).

The recursion \(P_0=1\), \(W_k=\Phi_m(mP_k)\),
\(P_{k+1}=P_kW_k\) has \(P_k=\prod_{i<k}W_i\) and pairwise
coprime \(W_i>1\), by the same constant-term calculation. It
therefore introduces a new prime at every step. Multiplying an input
also by a factorial \(b!\) excludes every prime at most \(b\).
These statements prove newness and, with such a factorial, strict
growth; (BP8)--(BP9) concern the entire family and do not assert
coverage by a particular single recursion. These are the mathematical
scopes of transcript lines 5438--6411 and 8031--8136.

\begin{theorem}[Short homogeneous \(\Phi_{12}\) factorization]
\label{thm:bp-short}
For every prime \(p\equiv1\pmod{12}\), there exist coprime
positive integers \(r,s\le\lfloor\sqrt p\rfloor\) such that
\[
 r^4-r^2s^2+s^4=ph,\qquad 1\le h<p.
 \tag{BP11}
\]
Every prime divisor of \(h\) is congruent to 1 modulo 12.
\end{theorem}
\begin{proof}
Choose \(w\) of order 12 modulo \(p\) and put
\(H=\lfloor\sqrt p\rfloor\). The map
\((i,j)\mapsto i-wj\) from \(\{0,\ldots,H\}^2\) into
\(\mathbb F_p\) has a repeated value because \((H+1)^2>p\).
The difference gives \((r_0,s_0)\ne(0,0)\),
\(|r_0|,|s_0|\le H\), \(r_0\equiv ws_0\pmod p\).
Neither coordinate vanishes: one zero would force the other to be
a multiple of \(p\) of absolute value less than \(p\).
Retain \(g=\gcd(|r_0|,|s_0|)\) and their two signs, and set
\(r=|r_0|/g,s=|s_0|/g\). These retained data recover the original
signed pair. Since \(g<p\), division by it is valid modulo \(p\).
Evenness in each coordinate gives
\(p\mid F(r,s):=r^4-r^2s^2+s^4\).

One has \(F=(r^2-s^2)^2+r^2s^2>0\). If \(r\ge s\),
then \(F=r^4-s^2(r^2-s^2)\le r^4\), and the opposite case is
symmetric. Thus \(0<F\le H^4<p^2\), giving (BP11).
If a prime \(q\mid F\) divided \(r\) or \(s\), it would
divide both, contrary to coprimality. A primitive pair gives
\(F\equiv1\pmod2\) and \(F\equiv1\pmod3\), so these two
primes are excluded. The unit \(rs^{-1}\) is a root of
\(\Phi_{12}\) modulo \(q\), and the order argument above gives
\(q\equiv1\pmod{12}\).
\end{proof}

This proves the assistant proposition at transcript lines 6612--6751
with the terminal alternative \(h=1\) retained. For comparison of
the cofactor and quadratic sign, there is an exact lift formula.
If \(F(r,s)=ph\), \(p\nmid rs\), then
\[
 \frac{F(r+pt,s)}p\equiv h+t\,2r(2r^2-s^2)\pmod p.
 \tag{BP12}
\]
The coefficient is a unit: \(2r^2=s^2\) would imply
\(F=3s^4/4=0\pmod p\), impossible. Hence unrestricted lifts
realize every quotient residue modulo \(p\), including either
nonzero square class. Primitivity can be retained by choosing \(t\)
to avoid the one forbidden class modulo each prime dividing \(s\),
using CRT together with the prescribed class modulo \(p\).
A positive sufficiently large representative gives a positive lift.
This exact statement does not retain the short-box bound.

\subsection{Necessary nonresidue factors and the exact character target}

\begin{lemma}[Divisor-of-a-sum obstruction]\label{lem:bp-square}
Let \(a,b,d>0\), \(\gcd(a,b)=1\), \(d\mid a+b\), and
\(d\equiv3\pmod4\). Then \(-d\) is not a square modulo \(4ab\).
\end{lemma}
\begin{proof}
Suppose \(t^2\equiv-d\pmod{4ab}\). The hypotheses give
\(\gcd(d,ab)=1\). For each odd prime \(\ell\mid ab\),
\(\left(\frac{-d}{\ell}\right)=1\). Quadratic reciprocity,
including its sign because \(d\equiv3\pmod4\), gives
\[
 \left(\frac\ell d\right)=\left(\frac{-d}\ell\right)=1.
\]
If \(2\mid ab\), reducing modulo 8 makes \(t\) odd and gives
\(d\equiv7\pmod8\), so \(\left(\frac2d\right)=1\).
Prime factorization and multiplicativity therefore imply
\(\left(\frac ad\right)=\left(\frac bd\right)=1\).
But \(a\equiv-b\pmod d\) gives
\(\left(\frac ad\right)=-\left(\frac bd\right)\), a contradiction.
\end{proof}

\begin{theorem}[A nonresidue prime in every normalized witness]
\label{thm:bp-nonresidue}
Every middle witness in (BP1) has a prime divisor \(\ell\mid AB\)
with \(\left(\frac\ell p\right)=-1\). Every exterior witness
has such a prime divisor of \(BC\).
\end{theorem}
\begin{proof}
For the middle channel, \(R\mid A+B\), \(\gcd(A,B)=1\),
and \(p\equiv-R\pmod{4AB}\). Lemma~\ref{lem:bp-square}
says \(p\) is not a square modulo \(4AB\). If every prime
\(\ell\mid AB\) were a residue modulo \(p\), reciprocity
would give \(\left(\frac p\ell\right)=1\) for each odd
such prime, because \(p\equiv1\pmod4\). A unit square root
modulo \(\ell\) lifts to every power: given \(x^2\equiv p\)
modulo \(\ell^e\), solve
\(2xt\equiv(p-x^2)/\ell^e\pmod\ell\) and replace
\(x\) by \(x+t\ell^e\). If \(2\mid AB\),
\(\left(\frac2p\right)=1\) and \(p\equiv1\pmod4\)
give \(p\equiv1\pmod8\). Every such unit is a square modulo
each 2-power: starting modulo 8, replacing an odd root \(x\)
modulo \(2^e\), \(e\ge3\), by \(x+t2^{e-1}\) toggles
the next square bit modulo \(2^{e+1}\). If \(AB\) is odd,
the needed modulo-4 root follows from \(p\equiv1\pmod4\).
CRT now makes \(p\) a square modulo \(4AB\), a contradiction.

For the exterior channel the tuple equation is
\[
 A(4BCD-1)=p(B+C).
\]
Since \(A<p\), set \(T=(B+C)/A\in\mathbb Z_{>0}\).
Then \(pT+1=4BCD\), \(T\equiv3\pmod4\), and \(T\mid B+C\).
A common prime divisor of \(B,C\) would divide \(A\), so
\(\gcd(B,C)=1\). Also \(p\nmid BC\): \(B<p\), and
\(p\mid C\) in the tuple equation would force \(p\mid A\).
If \(p\) were a square modulo \(4BC\), its inverse would be
one, contrary to Lemma~\ref{lem:bp-square} because
\(p^{-1}\equiv-T\pmod{4BC}\). The same prime-power and
reciprocity argument therefore supplies a nonresidue prime in \(BC\).
\end{proof}

For any fixed \(K\ge2\), put
\[
 m_K=\operatorname{lcm}\left(12n,8,
       \prod_{\substack{3\le\ell\le K\\\ell\ \text{prime}}}\ell\right).
\]
Every prime \(p\equiv1\pmod{m_K}\) makes all primes at most
\(K\) quadratic residues: use reciprocity for odd primes and the
modulo-8 supplementary law for 2. Theorem~\ref{thm:bp-nonresidue}
therefore forces a prime greater than \(K\) in the specified product
\(AB\) or \(BC\). In the middle channel, \(A\le2\) forces
that prime to divide \(B\), hence \(B>K\). Theorem
\ref{thm:bp-cyclotomic} produces infinitely many primes satisfying
these restrictions. This excludes fixed bounds on these pairs, not
all witnesses or the global two-layer assertion. These are the exact
quantifiers of transcript lines 6836--7265.

\begin{proposition}[A nonempty Jacobi-sign class and exact target counts]
\label{prop:bp-character}
Let \(p\ge13\) be prime and \(p\equiv1\pmod4\), and put
\(L=(p-1)/4\). There is a prime quadratic nonresidue
\(\ell\le L\). The integer
\(a=\ell\lceil(L+1)/\ell\rceil\) is a first-half shell denominator
containing a nonresidue prime factor. For this shell put
\[
 \tau=\prod_{r\mid a}(2v_r(a)+1),\qquad
 \sigma=\prod_{\substack{r\mid a\\(r/p)=1}}(2v_r(a)+1).
\]
Then the exact number of divisors \(u\mid a^2\) with
\(\left(\frac uR\right)=-1\) is
\((\tau-\sigma)/2\ge\tau/3>0\).

For every first-half shell, independently of this choice, let
\[
 W_a=\left\{\prod_{\ell\mid a}\ell^{e_\ell}\pmod R:
                   -v_\ell(a)\le e_\ell\le v_\ell(a)\right\}.
\]
Then \(W_a^{-1}=W_a\) and
\[
 E_a\cup M_a\ne\varnothing
 \quad\Longleftrightarrow\quad \{-1,-p\}\cap W_a\ne\varnothing.
 \tag{BP13}
\]
With \(\psi\) ranging over all characters of
\((\mathbb Z/R\mathbb Z)^\times\),
\[
 |E_a|+|M_a|=
 \frac1{\varphi(R)}\sum_\psi\psi(-1)(1+\psi(p))
       \prod_{\ell^v\parallel a}\sum_{e=-v}^v\psi(\ell)^e.
 \tag{BP14}
\]
\end{proposition}
\begin{proof}
If every integer 1 through \(L\) were a residue, their negatives
would also be residues because \(p\equiv1\pmod4\), exhausting
all \(2L\) nonzero residues. Every integer from \(L+1\) through
\(3L\) would then be a nonresidue. But \(2L=2\cdot L\) is a
product of two assumed residues and lies in that interval. The least
nonresidue is prime because otherwise all its prime factors are
smaller residues. The prescribed \(a\) satisfies
\(L+1\le a\le L+\ell\le2L\).

For an odd prime \(r\mid a\), reciprocity and
\(R\equiv-p\pmod r\), \(R\equiv3\pmod4\), give
\(\left(\frac rR\right)=\left(\frac rp\right)\): the two
reciprocity signs cancel. For \(r=2\), evenness of \(a\)
gives \(R\equiv-p\pmod8\) and the same identity by the
supplementary law. Summing the Jacobi symbols of divisors of \(a^2\)
therefore factors into geometric sums. A residue prime contributes
\(2v+1\), and a nonresidue prime contributes
\(\sum_{e=0}^{2v}(-1)^e=1\). The sum is \(\sigma\), the
total is \(\tau\), and the negative count follows. At least one
nonresidue factor gives \(\tau\ge3\sigma\).

Dividing \(u\mid a^2\) by \(a\) translates exponents by
\(-v_\ell(a)\), and inversion negates every centered exponent.
The middle target becomes \(-1\), while the exterior target
becomes \(-p^{-1}\); inversion changes it to \(-p\). This
proves (BP13). Orthogonality gives, for a unit \(t\),
\[
 \#\{u\mid a^2:u\equiv t\pmod R\}
 =\frac1{\varphi(R)}\sum_\psi\overline{\psi(t)}
       \prod_{\ell^v\parallel a}\sum_{e=0}^{2v}\psi(\ell)^e.
\]
Shift each exponent by \(-v\), retaining the factor \(\psi(a)\).
At \(t=-a\) its prefactor is \(\psi(-1)\); at
\(t=-4^{-1}\) it is \(\psi(-1)\psi(4a)=\psi(-1)\psi(p)\).
Adding gives (BP14). No individual summand is asserted nonnegative;
the total is an integer target count.
\end{proof}

There is a further exact comparison with the divisor-pair statistic
used in the dialogue. Let
\[
 T_{R,S}=\#\{(m,r):m,r\mid S,\ R\mid m+r\}.
\]
Character orthogonality gives
\[
 T_{R,S}=\frac1{\varphi(R)}\sum_\psi\psi(-1)
                  \left|\sum_{d\mid S}\psi(d)\right|^2.
 \tag{BP15}
\]
Indeed expansion has the target \(-mr^{-1}=1\) in the unit group.
Reduction by \(g=\gcd(m,r)\) maps these pairs onto the primitive
pairs \((m_0,r_0)\) with \(m_0,r_0\mid S\),
\(\gcd(m_0,r_0)=1\), \(R\mid m_0+r_0\). Its full fibre is
\[
 \{(g m_0,g r_0):g\mid S/(m_0r_0)\}.
 \tag{BP16}
\]
To prove this, division by \(g\mid S\) preserves the congruence
because \(g\) is a unit modulo \(R\). Conversely both scaled
entries divide \(S\) exactly when
\(g\mid\gcd(S/m_0,S/r_0)=S/(m_0r_0)\). The primitive pair
reconstructs the ordered solution
\[
 k=(m_0+r_0)/R,\qquad
 (a,y,z)=(a,(S/r_0)k,(S/m_0)k).
\]
Its reciprocal sum is \(1/a+R/S=4/p\). The inverse primitive
pair is the reduced positive ratio \((Ry-S)/S=y/z\); prime
valuations in \((Ry-S)(Rz-S)=S^2\) prove both its entries
divide \(S\). Thus \(T>0\) is exactly shell occupancy, while
its multiplicities are the divisor counts of (BP16). A canonical
hit supplies the distinct pair \((A,p^{1-\varepsilon}B)\),
so \(T\ge2\). This proves, with its full multiplicity map,
the relation in transcript lines 3594--3744 and 7843--7887.

\subsection{Five exact regressions and their scopes}

The following finite certificates are the five stronger-claim
counterexamples of transcript lines 8330--8786. They were reproduced
by a separate standard-library integer calculation. Their exhaustiveness
concerns the displayed finite divisor sets or coordinate boxes, not all
primes or all Erd\H{o}s--Straus candidates.

\paragraph{1. A selected nonresidue shell can be empty.}
Take \(p=3361\), \(L=840\). The primes 2,3,5,7 are residues
modulo \(p\): use \(p\equiv1\pmod8\), \(p\equiv1\pmod{12}\),
and reciprocity with \(p\equiv1\pmod5\), \(p\equiv1\pmod7\).
Thus every integer 1 through 10 is a residue. On the other hand
\(3361\equiv6\pmod{11}\), and the nonzero squares modulo 11
are \(\{1,3,4,5,9\}\), so 11 is the least nonresidue.
Proposition~\ref{prop:bp-character} selects
\[
 a=11\lceil841/11\rceil=847=7\cdot11^2,\quad R=27,\quad j=6.
\]
Every divisor of \(a^2\) is \(7^i11^e\), \(0\le i\le2\),
\(0\le e\le4\). Their residues are exactly
\[
 \{1,2,7,8,10,11,13,14,16,19,22,23,26\}\pmod{27}.
\]
They omit \(-4^{-1}=20\) and \(-a=17\), proving
\(E_6=M_6=\varnothing\). The six negative-Jacobi divisors are
\(11,77,539,1331,9317,65219\). The principal contribution
to (BP14) is \(2\cdot15/18=5/3\), while its complete sum is
zero; the other characters contribute \(-5/3\).
This is not an ES counterexample: the first shell has \(a=841\),
\(R=3\), \(u=29\mid841^2\), and \(3\mid4\cdot29+1\).
Its tuple is \((A,B,C,D,\varepsilon)=(29,1,1130,29,0)\), giving
\[
 c_H=79027200,\qquad
 \frac4{3361}=\frac1{841}+\frac1{950330}+\frac1{110139970}.
\]

\paragraph{2. Nonempty local fibres need not intersect in the divisor box.}
At \(p=37,a=18,R=35\), the complete divisor list is
\[
 \operatorname{Div}(324)=
 \{1,2,3,4,6,9,12,18,27,36,54,81,108,162,324\}.
\]
Those satisfying \(5\mid4u+1\) are \(1,6,36,81\); those
satisfying \(7\mid4u+1\) are \(12,54\). Both lists are nonempty
and their intersection is empty. CRT gives \(u\equiv26\pmod{35}\),
but no member of the original complete list belongs to that class.
This refutes the exchange of \(\forall m\exists u\) with
\(\exists u\forall m\), while preserving Theorem~\ref{thm:bp-crt}.

\paragraph{3. The short cofactor may always be one at a nonbase prime.}
The complete positive primitive boxes in Theorem~\ref{thm:bp-short}
give the following table:
\[
\begin{array}{c|c|c|c}
 p&H&\text{all primitive pairs}&h\\\hline
 61&7&(2,3),(3,2)&1\\
 73&8&(1,3),(3,1)&1\\
 193&13&(3,4),(4,3)&1\\
 2521&50&(1,26),(26,1)&181
\end{array}
\]
Here is a finite exact verification of completeness. The roots of
\(\Phi_{12}\) in these four prime fields are, respectively,
\[
 \{21,29,32,40\},\quad \{3,24,49,70\},\quad
 \{49,63,130,144\},\quad \{26,97,2424,2495\}.
\]
Substitution gives four distinct roots in each case, so the polynomial
degree proves the lists complete. For each root \(w\), the possible
positive coordinates in the box are exactly
\[
 r=ws-pk,\quad
 \max\left(1,\left\lceil\frac{pk+1}{w}\right\rceil\right)
 \le s\le
 \min\left(H,\left\lfloor\frac{pk+H}{w}\right\rfloor\right),
 \quad0\le k\le\lfloor wH/p\rfloor.
 \tag{BP17}
\]
Evaluating these integer intervals gives, before removing nonprimitive
pairs, respectively
\[
\begin{array}{c|l}
61&(2,3),(3,2),(4,6),(6,4)\\
73&(1,3),(2,6),(3,1),(6,2)\\
193&(3,4),(4,3),(6,8),(8,6),(9,12),(12,9)\\
2521&(1,26),(26,1).
\end{array}
\]
The gcd test and direct quartic evaluation give the preceding table.
For example \(2^4-2^2 3^2+3^4=61\). Hence \(h>1\) cannot
be demanded at every prime greater than 13 in this prescribed box.

\paragraph{4. A smaller cyclotomic prime need not be a nonresidue.}
The last row of that complete table gives
\[
 1-26^2+26^4=456301=2521\cdot181,
 \qquad 88^2=3\cdot2521+181.
\]
Thus \(181<2521\), \(181\equiv1\pmod{12}\), but
\(\left(\frac{181}{2521}\right)=1\). Both permitted pairs
give that cofactor. The actual middle tuple
\((A,B,C,D,\varepsilon)=(2,159,7,2,1)\) has
\(AB=2\cdot3\cdot53\), with 53 a nonresidue modulo 2521:
reciprocity gives
\[
 \left(\frac{53}{2521}\right)
 =\left(\frac{30}{53}\right)
 =\left(\frac2{53}\right)\left(\frac3{53}\right)
  \left(\frac5{53}\right)=(-1)(-1)(-1)=-1.
\]
Its code and original denominators are
\[
 c_H=1600517,\qquad
 \frac4{2521}=\frac1{636}+\frac1{70588}+\frac1{5611746}.
\]
Substitution gives \(4ABCD=17808=A+B+pC\), and all the code
bits equal one. Thus the cofactor and witness factor have their
respective exact roles.

\paragraph{5. The degree-four cofactor bound does not persist for
\(\Phi_{24}\).}
For \(p=73\), the eight roots of \(\Phi_{24}(X)=X^8-X^4+1\)
modulo 73 are \(7,17,21,30,43,52,56,66\). Substitution and degree
prove completeness. Formula (BP17), now with \(H=8\) and this root
list, gives exactly
\((1,7),(7,1),(4,5),(5,4)\). They are all primitive, and
\[
 F_{24}(1,7)=5762401=73\cdot78937,\quad
 F_{24}(4,5)=296161=73\cdot4057.
\]
The smallest quotient is \(4057>73\), disproving the unchanged
requirement \(h<p\). The valid bound for this degree-eight trinomial
is \(F_{24}(r,s)\le\max(r,s)^8<p^4\), hence \(h<p^3\).
The cyclotomic infinitude theorem at modulus 24 is unaffected.

The primality checks needed in these finite certificates are elementary
trial divisions by all primes at most the respective square roots.
For 61 and 73 those primes are \(2,3,5,7\); for 193 they are
\(2,3,5,7,11,13\); for 2521 they are
\(2,3,5,7,11,13,17,19,23,29,31,37,41,43,47\); for 3361
add 53 to that last list. None divides the corresponding number.
For the auxiliary primes 97,181,197 the same test uses respectively
the lists through 7,13,13. The complete finite divisor sets, root
lists, quotient values and reciprocal identities above are exact
certificates for their stated counterexamples. They do not establish
the universal positivity of (BP3) or (BP9).

% Reconstructed proof module, 12 September 2026.
% This is newly authored mathematics, not a recovered src/energy.tex download.
% Required packages: amsmath, amssymb, amsthm; the standalone wrapper supplies them.
\section{Exact character-resolved collision minima and integral ES return}
\label{sec:reconstructed-character-energy}

\paragraph{Provenance and scope.}
This module reconstructs the character-resolved energy result stated in the
supplied joint ES discussion, \texttt{pasted-text.txt}, in its final
``Today 8:19 AM'' entry, subsection ``JT's contribution supplies an exact
refinement of the positivity test.'' The discussion credits the quadratic
character refinement to its integration of JT's bounded-divisor synthesis.
The original \texttt{src/energy.tex} was not available: the downloaded loose
\texttt{workbench.tex} is a wrapper referring to that missing file.
The shell and ordered channel conventions are those in the supplied earlier
\texttt{es\_dual\_frontier/workbench.tex} and the definition block in
\texttt{dual\_es.py}; their required identities are proved here in full.
The independently authored checker imports no downloaded code.
The result is a sufficient test on a specified full exponent box. It proves
no universal statement that every prime has a shell passing this test.

\subsection{The complete shell, exponent box, and ordered return}
Fix a prime $p=4m+1$, where $m\geq1$, and an integer
\[
 m+1\leq a\leq3m,\qquad R=4a-p.
\]
Then $3\leq R\leq8m-1=2p-3$, $R\equiv3\pmod4$, and $0<a<p$.
Since $p$ is prime, $\gcd(a,p)=1$; consequently
\[
 \gcd(a,R)=\gcd(a,p)=1,\qquad
 \gcd(p,R)=\gcd(p,4a)=1.
\]
Write the complete prime factorization as $a=\prod_{q\mid a}q^{e_q}$.
All these primes are units modulo $R$. Define
\[
 B_a=\prod_{q\mid a}\bigl([-e_q,e_q]\cap\mathbb Z\bigr),
 \quad v(\beta)=\prod_{q\mid a}q^{\beta_q}\in
 G=(\mathbb Z/R\mathbb Z)^\times,
 \quad u(\beta)=\prod_{q\mid a}q^{e_q+\beta_q}.
\]
Negative powers in $v$ mean multiplicative inverses in $G$.
Unique prime factorization proves that $\beta\mapsto u(\beta)$ is a bijection
from $B_a$ onto the positive divisors of $a^2$: its inverse is
$\beta_q=v_q(u)-e_q$. Multiplying the residue products gives
\[
 v(\beta)=u(\beta)a^{-1}\pmod R,
 \qquad u(-\beta)=a^2/u(\beta),
 \qquad v(-\beta)=v(\beta)^{-1}.
\]
In particular no unbounded subgroup or shorter exponent interval has
replaced $B_a$.

For any positive divisor $u\mid a^2$, put $d=a^2/u$ and retain the two
ordered rational triples
\begin{align*}
 Q_E(p,a,u)&=\left(a,\frac{pa+d}{R},\frac{pa+p^2u}{R}\right),\\
 Q_M(p,a,u)&=\left(a,\frac{p(a+d)}{R},\frac{p(a+u)}{R}\right).
\end{align*}
Their first denominator remains distinguished; the last two entries are
not sorted. In either case, if these entries are denoted $Y,Z$, then
\[
 (RY-pa)(RZ-pa)=p^2a^2.
\]
Indeed the two factors are $(d,p^2u)$ in channel $E$ and $(pd,pu)$ in
channel $M$. Expanding and cancelling the identical constant term gives
$R^2YZ-paR(Y+Z)=0$. Since $R,p,a,Y,Z$ are positive, division gives
\[
 \frac1Y+\frac1Z=\frac{R}{pa},\qquad
 \frac1a+\frac1Y+\frac1Z=\frac{p+R}{pa}=\frac4p.
\]
Both triples therefore always give positive rational ES decompositions.

The exact integral gates follow without discarding the modulus $R$.
For $E$, multiplication of the first numerator by the unit $u$ yields
\[
 u(pa+d)\equiv a^2(4u+1)\pmod R,
 \qquad pa+p^2u\equiv4a^2(4u+1)\pmod R.
\]
For $M$ the exact numerator identities are
\[
 u\,p(a+d)=pa(u+a),\qquad p(a+u)=p(u+a).
\]
Since every multiplier used is a unit modulo $R$, the two reduced
denominators in each channel coincide and are respectively
\[
 D_E=\frac{R}{\gcd(R,4u+1)},\qquad
 D_M=\frac{R}{\gcd(R,u+a)}.
\]
Thus $Q_E$ is integral exactly when $R\mid4u+1$, and $Q_M$ exactly when
$R\mid u+a$. This proves all assertions about integrality, not merely
one implication.

Since $u=av(\beta)$ modulo $R$ and $4a=p$ modulo $R$, these tests become
\[
 M:\ v(\beta)=-1,
 \qquad E:\ v(\beta)=-p^{-1}.
\]
The complete box is stable under $\beta\mapsto-\beta$, and
$(-p)^{-1}=-p^{-1}$, so a vector with residue $-p$ also supplies an
exterior witness by applying this involution. Consequently simultaneous
failure of both channels in this shell is equivalent to absence of all
the residues in the deduplicated, inversion-stable set
\[
 T=\{-1,-p,-p^{-1}\}\subseteq G.
\]
The set may have fewer than three elements, and may contain the identity;
both possibilities are retained below.

For the additional integral marking put
\[
 g=\gcd(a,u),\quad h=g^2/u,\quad r=u/g,\quad s=a/g.
\]
If $e=v_q(a)$ and $f=v_q(u)$, then $0\leq f\leq2e$ and the exponents
of $h,r,s$ are $2\min(e,f)-f$, $f-\min(e,f)$ and
$e-\min(e,f)$. The first is $f$ when $f\leq e$ and $2e-f$ otherwise,
so all are nonnegative. At least one of the last two is zero, proving
$\gcd(r,s)=1$. Direct substitution now proves
\[
 a=hrs,\quad u=hr^2,\quad d=hs^2,
\]
and the ordered triples are
\[
 Q_E=(hrs,hs\kappa,phr\kappa),\quad\kappa=(pr+s)/R,
 \qquad
 Q_M=(hrs,phs\lambda,phr\lambda),\quad\lambda=(r+s)/R.
\]
Since $g$ is a unit modulo $R$, the middle gate is equivalent to
$R\mid r+s$. The exterior numerator satisfies
$pa+d=hs(pr+s)$, where $hs$ is a unit modulo $R$, so the exterior
gate is equivalent to $R\mid pr+s$. Hence these quotients are positive
integers whenever the corresponding target occurs. Every marking returns
to the original objects via $a=hrs$, $u=hr^2$ and
$\beta_q=v_q(u)-e_q$; moreover
\[
 u=\frac{a^2}{RY-pa}\quad(E),\qquad
 u=\frac{pa^2}{RY-pa}\quad(M).
\]
Thus the channel and ordered denominators retain an explicit inverse.

For completeness, the full interval of shells suffices for the ES
existence question at this $p$. Given any positive integral solution,
move a smallest denominator to the distinguished position $a$, retaining
that permutation. The inequalities
$4/p>1/a$ and $4/p\leq3/a$ give $p/4<a\leq3p/4$, hence
$m+1\leq a\leq3m$. Its other two denominators obey the residual product
identity above, and
$RY-pa=paY/Z>0$, $RZ-pa=paZ/Y>0$. Since $p\nmid a$, the first residual
is uniquely $p^i v$ with $i\in\{0,1,2\}$ and $v\mid a^2$.
For $i=0$, take $E,u=a^2/v$ in the existing order. For $i=1$, take
$M,u=a^2/v$ in the existing order. For $i=2$, take $E,u=v$ and retain
the swap of the last two entries. The displayed residuals verify each
case. Conversely, an integral candidate already solves ES. This proves
the exact equivalence with a hit in the full atlas; it does not assert
that such a hit always exists.

\subsection{Multiplicities, parity, and character cells}
Define
\[
 \mu(g)=\#\{\beta\in B_a:v(\beta)=g\},\quad
 N=\sum_{g\in G}\mu(g)=\prod_{q\mid a}(2e_q+1),\quad
 E=\sum_{g\in G}\mu(g)^2.
\]
The involution $\beta\mapsto-\beta$ bijects the fibres of $g$ and
$g^{-1}$, proving $\mu(g)=\mu(g^{-1})$. Its only fixed vector in the
integer box is zero. The fibre over $1$ contains that vector and otherwise
splits into two-element orbits; thus $\mu(1)$ is odd and at least one.
If $g^2=1$ but $g\ne1$, its fibre is stable under this involution and
contains no fixed vector, so $\mu(g)$ is even.

The collision count also has the independent exact expression
\[
 E=\sum_{\substack{-2e_q\leq\delta_q\leq2e_q\ (q\mid a)\\
                  \prod q^{\delta_q}=1\ (\bmod R)}}
       \prod_{q\mid a}(2e_q+1-|\delta_q|).
\]
Indeed $E$ counts ordered pairs $(\beta,\gamma)$ with identical residue.
Their coordinate difference $\delta=\beta-\gamma$ has residue one.
For a single interval $[-e,e]$, the number of ordered pairs of difference
$\delta$ is $2e+1-|\delta|$ for $|\delta|\leq2e$, and zero otherwise:
the first entry lies in $[-e,e]\cap[-e+\delta,e+\delta]$.
Taking the product of these counts proves the formula.

Let $V$ be a finite elementary abelian $2$-group, written
multiplicatively, and let $\pi:G\to V$ be a homomorphism; surjectivity
is not required. For $v\in V$ put
\[
 C_v=\pi^{-1}(v),\qquad
 N_v=\sum_{g\in C_v}\mu(g),\qquad
 E_v=\sum_{g\in C_v}\mu(g)^2.
\]
As $v^{-1}=v$, each cell is stable under inversion. Its mass is even
except in the cell containing $1$, where its mass is odd: inverse pairs
contribute twice their common multiplicity, nonidentity involutions have
even multiplicity, and the identity contributes odd multiplicity.
Empty cells have mass and energy zero.

For each odd prime $\ell\mid R$, the prime-local quadratic character
$\chi_\ell(g)$ is $1$ when $g$ is a nonzero square modulo $\ell$ and
$-1$ otherwise. This is a homomorphism: in the multiplicative group of
the field $\mathbb Z/\ell\mathbb Z$, the square map is a homomorphism
whose kernel is exactly $\{1,-1\}$, since
$(x-1)(x+1)=0$ implies $x=1$ or $x=-1$ in a field. Every fibre has
two elements, so the square subgroup has index two. The quotient by
that subgroup has two elements and therefore supplies precisely this
sign character. Reduction modulo $\ell$ is a group homomorphism on $G$,
which proves the assertion for $\chi_\ell$ on $G$ as well.

Writing $R=\prod_{\ell\mid R}\ell^{f_\ell}$, define
\[
 \pi_{\rm quad}(g)=(\chi_\ell(g))_{\ell\mid R},\qquad
 \pi_{\rm Jac}(g)=\prod_{\ell\mid R}\chi_\ell(g)^{f_\ell}.
\]
Both are homomorphisms. The second is obtained from the first by the
explicit homomorphism $(\epsilon_\ell)\mapsto
\prod\epsilon_\ell^{f_\ell}$; factors with even $f_\ell$ disappear only
from this product. For example, modulo $9$ the residue $2$ is a
non-square modulo $3$, so its prime-local character is $-1$, whereas
its Jacobi character is $(-1)^2=1$. Thus these two partitions need
not coincide. The trivial homomorphism is the unpartitioned case.

\subsection{The exact minimum, including all impossible cases}
For a fixed cell let $A_v=C_v\setminus T$. Since both $C_v$ and $T$
are inversion-stable, so is $A_v$. Its slots are the unordered inverse
pairs $\{g,g^{-1}\}$ with $g\ne g^{-1}$, the individual nonidentity
involutions $g^2=1$, $g\ne1$, and the identity if it is present.
Write $b_v=1$ when $1\in C_v$ and $b_v=0$ otherwise.

For a nonnegative integer mass $n$, define $L_v(n)$ as the minimum of
$\sum_{g\in C_v}\nu(g)^2$ over all maps
$\nu:C_v\to\mathbb Z_{\geq0}$ satisfying
\[
 \sum_{g\in C_v}\nu(g)=n,\quad
 \nu(g)=\nu(g^{-1}),\quad
 \nu|_{T\cap C_v}=0,
\]
and additionally $\nu(1)$ odd when $1\in C_v$, and $\nu(g)$ even
for every nonidentity involution in the cell. If there is no such map,
set $L_v(n)=+\infty$.

Here is a complete evaluation. If $1\in T\cap C_v$, the minimum is
$+\infty$ for every $n$, since an odd nonnegative identity value cannot
vanish. Otherwise the identity slot exists precisely when $b_v=1$.
If $n<b_v$ or $n-b_v$ is odd, the minimum is $+\infty$.
In all remaining cases put $K=(n-b_v)/2$. Write $P_v$ for the set
of inverse-pair slots and $J_v$ for the nonidentity involution slots.
Use nonnegative integer variables $k_P$ for $P\in P_v$, $z_j$ for
$j\in J_v$, and $k_0$ only if $b_v=1$. Then
\[
 L_v(n)=\min_{\sum_P k_P+\sum_j z_j+b_vk_0=K}
 \left(2\sum_P k_P^2+4\sum_j z_j^2+
 \begin{cases}(1+2k_0)^2,&b_v=1,\\0,&b_v=0.\end{cases}\right),
\]
where the absent variable is omitted from both sums. If $K>0$ and
there are no slots, the minimum is $+\infty$; if $K=0$ it is $b_v$.
These are all the impossible cases: when at least one slot exists, it
can hold every remaining pair of mass units. For a pair slot the two
values are $(k_P,k_P)$, for a nonidentity involution its value is
$2z_j$, and for the identity its value is $1+2k_0$. These assignments
give a bijection between feasible slot allocations and feasible maps
$\nu$, proving the optimization formula. There are finitely many
allocations for fixed $K$, so every finite minimum is attained.

Start the identity at mass one and energy one when $b_v=1$ and start
every other slot at zero. Incrementing a slot with current index $k$
by one adds two units of mass. Its exact energy increment is respectively
\[
 \begin{array}{c|c|c}
 \text{slot}&\text{current values}&\text{incremental energy}\\\hline
 \{g,g^{-1}\},\ g\ne g^{-1}&(k,k)&2(k+1)^2-2k^2=4k+2\\
 g^2=1,\ g\ne1&2k&(2k+2)^2-(2k)^2=8k+4\\
 g=1&1+2k&(3+2k)^2-(1+2k)^2=8k+8
 \end{array}
\]
In each row $k$ starts at zero. Thus the marginal sequences are
$2,6,10,\ldots$, $4,12,20,\ldots$, and $8,16,24,\ldots$.
Selecting the $K$ least costs from the union of these increasing
sequences proves the minimum exactly. To justify feasibility, every
predecessor in a sequence is strictly smaller than its successor,
so selecting a successor among the $K$ least costs selects its whole
prefix. Ties between different sequences can be resolved in any fixed
order. To justify optimality, any feasible allocation selects exactly
$K$ costs from this union, hence its sum is at least the sum of the
$K$ least ones. Equivalently, replacing a selected unnecessarily large
cost by the least unselected available cost improves the allocation
until the same minimum is reached. A heap holding the first unused
cost of each sequence implements this finite procedure.

\paragraph{Forcing theorem.}
For the actual full-box multiplicities and every cell $v$,
\[
 \mu|_{T\cap C_v}=0\quad\Longrightarrow\quad E_v\geq L_v(N_v).
\]
Indeed all other constraints defining $L_v$ were proved for $\mu$,
so this restriction is a feasible map if its target fibres vanish.
Taking the contrapositive proves
\[
 E_v<L_v(N_v)\quad\Longrightarrow\quad
 \text{there is }\beta\in B_a\text{ with }v(\beta)\in T\cap C_v.
\]
The preceding explicit map then produces a positive integral ES
witness. In particular, if $T\cap C_v=\{-1\}$, the witness is in
the middle channel. A finite energy is smaller than $+\infty$,
so incompatible cells are included without an extra convention.
If $1\in T$, the zero vector itself is already a target and is
detected by this rule. If $T\cap C_v$ is empty, the actual restriction
is feasible and this cell can never produce a false certificate.

Every finite $L_v(n)$ is attained by a distribution satisfying exactly
the specified mass, symmetry, parity, and exclusions. Hence it is the
largest universal lower bound on energy from those constraints alone.
This is the precise optimality assertion. It does not state that every
energy above the minimum is attainable, or that every attaining abstract
distribution arises from the arithmetic exponent box.

\subsection{Exact refinement and preservation of certificates}
Let $\pi_2:G\to V_2$ and $\pi_1:G\to V_1$ be such homomorphisms,
and suppose $\pi_1=\rho\circ\pi_2$ for a homomorphism
$\rho:V_2\to V_1$. A coarse cell $C_v$ is the disjoint union of
the fine cells $C_w$ with $\rho(w)=v$. For any prescribed fine
masses $n_w$, with $n=\sum_{\rho(w)=v}n_w$, the exact inequality is
\[
 L^{(1)}_v(n)\leq\sum_{\rho(w)=v} L^{(2)}_w(n_w).
\]
All terms take values in $\mathbb Z_{\geq0}\cup\{+\infty\}$.
If some right-hand term is infinite the inequality is automatic.
Otherwise choose a minimizer on each fine cell and join these maps.
Each inversion orbit lies entirely in a fine cell, so the joined map
obeys every coarse constraint, with mass $n$ and energy equal to the
sum on the right. This proves the inequality. Restricting any coarse
feasible map to its fine cells proves the stronger exact identity
\[
 L^{(1)}_v(n)=
 \min_{\substack{n_w\in\mathbb Z_{\geq0}\\
                  \sum_{\rho(w)=v}n_w=n}}
       \sum_{\rho(w)=v}L^{(2)}_w(n_w).
\]
If the coarse minimum is finite, its minimizer supplies the masses
and gives the reverse inequality. If it is infinite, no combination
of finite fine minimizers can exist, since their join would be a
coarse feasible map. This also proves equality in the infinite case.

For actual masses and energies,
$N^{(1)}_v=\sum N^{(2)}_w$ and $E^{(1)}_v=\sum E^{(2)}_w$ by disjointness.
If every fine cell satisfies $E^{(2)}_w\geq L^{(2)}_w(N^{(2)}_w)$,
their sum and the inequality above give
$E^{(1)}_v\geq L^{(1)}_v(N^{(1)}_v)$. The contrapositive proves that
every coarse forcing certificate yields a forcing certificate in at
least one of its fine cells. This applies successively to the trivial,
Jacobi, and all-prime-local quadratic partitions, with the explicit
product map already proved above.

\subsection{The complete numerical witness at p = 1201}
Take $p=1201$, $a=310=2\cdot5\cdot31$, $R=39=3\cdot13$.
To certify primality explicitly, $34^2<1201<35^2$ and the residues
of $1201$ modulo the primes $2,3,5,7,11,13,17,19,23,29,31$ are
respectively $1,1,1,4,2,5,11,4,5,12,23$. None is zero, and every
composite integer has a prime divisor at most its square root, as
follows by writing it as the product of two positive factors larger
than one and taking a prime factor of the smaller one. These are
all primes at most 34, so $1201$ is prime.
There are exactly $3^3=27$ exponent vectors with each coordinate
in $\{-1,0,1\}$. In the residue ring, $p=31$ and $31^{-1}=34$,
since $31\cdot34=27\cdot39+1$. Hence $T=\{5,8,38\}$.
The complete histogram on the 24 units is
\[
\begin{array}{c|rrrrrrrrrrrr}
g&1&2&4&5&7&8&10&11&14&16&17&19\\\hline
\mu(g)&1&1&1&1&2&1&1&0&2&1&1&2\\[2pt]
g&20&22&23&25&28&29&31&32&34&35&37&38\\\hline
\mu(g)&1&1&1&0&2&1&1&0&1&1&2&2
\end{array}
\]
This table can be checked by multiplying all 27 combinations of
$(2^{-1},1,2)=(20,1,2)$, $(5^{-1},1,5)=(8,1,5)$ and
$(31^{-1},1,31)=(34,1,31)$ modulo $39$. Its six entries of value
two and fifteen entries of value one give $N=27$ and $E=6\cdot4+15=39$.

Deleting $T$ leaves nine inverse pairs,
\[
 \{2,20\},\{4,10\},\{7,28\},\{11,32\},\{16,22\},
 \{17,23\},\{19,37\},\{29,35\},\{31,34\},
\]
the two nonidentity involutions $14,25$, and the identity. The
unpartitioned mass budget is $(27-1)/2=13$. The 13 least marginal
costs are nine copies of $2$, two copies of $4$, and two copies of
$6$. Thus
\[
 L_{\mathrm{unpartitioned}}(27)=1+9\cdot2+2\cdot4+2\cdot6=39.
\]
Equality $E=39$ gives no strict forcing certificate at this level.

The negative Jacobi cell is exactly
\[
 C_- =\{7,14,17,19,23,28,29,31,34,35,37,38\}.
\]
For a direct sign check the nonzero squares modulo $3$ are $\{1\}$,
and those modulo $13$ are $\{1,3,4,9,10,12\}$; taking the product
of these two signs gives exactly the displayed cell. Its table entries
give $N_-=18$, $E_-=30$, and $T\cap C_-=\{38\}$.
Removing this target leaves five pairs
\[
 \{7,28\},\{17,23\},\{19,37\},\{29,35\},\{31,34\}
\]
and the nonidentity involution $14$. There is no identity slot.
The budget is nine, whose nine least costs are five copies of $2$,
one copy of $4$, and three copies of $6$. Therefore
\[
 L_-(18)=5\cdot2+4+3\cdot6=32
 =\min_{n_1+\cdots+n_5+z=9}
       \left(2\sum_{i=1}^5 n_i^2+4z^2\right).
\]
All variables in the minimum are nonnegative integers. The allocation
$(n_1,\ldots,n_5,z)=(2,2,2,1,1,1)$ attains 32, and the marginal
argument proves that no smaller value occurs. Hence $30<32$ forces
the middle target in this specific cell.

The finer labels are ordered by the primes $(3,13)$; their complete
mass, energy, minimum and target data are
\[
\begin{array}{c|r|r|r|l}
v&N_v&E_v&L_v(N_v)&T\cap C_v\\\hline
(-1,-1)&4&4&4&\{5,8\}\\
(-1,+1)&8&12&14&\{38\}\\
(+1,-1)&10&18&18&\varnothing\\
(+1,+1)&5&5&5&\varnothing
\end{array}
\]
These values follow by grouping the full histogram using the explicit
square lists. The allowed slot counts are respectively two pairs;
two pairs and the involution $14$; three pairs; and two pairs, the
identity, and the involution $25$. The marginal formula gives the
four displayed minima, including $2+2+4+6=14$ for $(-1,+1)$.

An actual middle fibre vector is $\beta=(0,-1,-1)$ in the prime
order $(2,5,31)$. Its residue is $8\cdot34=272=6\cdot39+38$, and
its original divisor is $u=2^{1+0}5^{1-1}31^{1-1}=2$.
The complete marking is therefore
\[
 (p,R,u,h,r,s,\lambda)=(1201,39,2,2,1,155,4),
\]
because $\gcd(310,2)=2$ and $(1+155)/39=4$. Substitution in the
ordered middle return gives
\[
 Q_M=(310,\ 1201\cdot2\cdot155\cdot4,\ 1201\cdot2\cdot1\cdot4)
     =(310,1489240,9608).
\]
Every number is a positive integer. The earlier proved reciprocal
identity now verifies the ES equation exactly, and the retained
slope satisfies $r/s=1/155<1/8$. This is a stronger forcing
certificate for this witness, with its entire divisor vector retained.

\subsection{Finite verification and the boundary of the result}
The independent standard-library checker compares the marginal algorithm
with an independent dynamic program
that uses the occupancy squares directly. It includes all masses
$0\leq n\leq19$, zero through three pair slots, zero through three
nonidentity involution slots, presence or absence of the identity,
and the forbidden-identity case. It also compares the actual cells
of the $1201,310$ example with that dynamic program. The checker
uses explicit exceptions, so optimized Python retains every check.
For its arithmetic cases it enumerates the full box and independently
enumerates divisors of $a^2$ by trial divisibility and complementary
divisors. The signed-box symmetries, parity, target exclusions,
refinement inequalities, and extracted rational identities are checked
exactly. The weighted autocorrelation is independently checked on
the listed examples, including shells with $R\geq p$.

The scan uses exactly the primes $p\equiv1\pmod{12}$ with $p\leq1500$
and the first-half shells $m+1\leq a\leq2m$, equivalently $R<p$.
It contains 54 primes and 9531 shells and gives respectively
214 unpartitioned, 231 Jacobi, and 240 fully quadratic certifications.
These are shell counts. The general theorem retains the full shell
range through $3m$; this finite scan does not replace that range.
The empty shell $p=37,a=18,R=35$ is retained and is not certified by
any partition. Ordinary and optimized execution reproduce the same
mathematical JSON outputs, with execution details recorded separately.
These finite results and the proved sufficient criterion do not
establish that every prime has a certified shell, and a shell that
does not pass this strict inequality may still have a genuine target.

\section{The original gate in an integral three-coordinate system}
This section reconstructs the final supplied discussion of 12 September,
08:19, from its complete text. The downloaded \texttt{workbench.tex} is only
a wrapper: its five mathematical input files were not recovered. The formulas
below are therefore explicitly a reconstruction, not an assertion that those
missing source bytes or their execution receipts have been recovered.
The comparison to JT's quartic calculation is a contribution of the supplied
joint investigation; the proofs here check the stated maps independently.

Fix an odd positive integer $R$ and an integer gate $g$. All four marked
coordinates $(g,x,y,z)$ are retained, with $g$ fixed by every map. On
$(\mathbb Z/R\mathbb Z)^3$ define
\begin{align*}
 a_1^g(x,y,z)&=(6g+y,-6g-x-y,z-2g+x),\\
 a_2^g(x,y,z)&=(-y,x-6g,z+3g+y),\\
 a_\infty^g(x,y,z)&=(x,y+x,z-g).
\end{align*}
These are the affine slices of the integral matrices in the received
\nolinkurl{es_s6_counterfactual/src/core.tex}, equation (1), whose source is
the $(3,4,\infty)$ period manuscript circulated by Levent Alp\"oge.
Only the displayed integral representation is used; no assertion about the
global diffeomorphism type of its proposed completed threefold is needed.
Direct composition gives
$(a_1^g)^3=(a_2^g)^4=a_1^ga_2^ga_\infty^g=1$.
The inverses are $(a_1^g)^2,(a_2^g)^3$ and
$(x,y,z)\mapsto(x,y-x,z+g)$ respectively.

\begin{proposition}[Integral cyclic coordinates]
For each retained $g$, the map
\[
 t_0=-x-y-z-2g,\qquad t_1=-x-z+2g,\qquad t_2=-z
\]
is an affine isomorphism over $\mathbb Z$, hence modulo every $R$. Its inverse is
\[
 x=t_2-t_1+2g,\qquad y=-t_0+t_1-4g,\qquad z=-t_2.
\]
In these coordinates the three maps become, respectively,
\begin{align*}
 c(t)&=(t_2,t_0,t_1),\\
 b_g(t)&=(t_1-g,t_2-g,t_0-t_1+t_2+g),\\
 u_g(t)&=(t_0+t_1-t_2-g,t_1+g,t_2+g).
\end{align*}
\end{proposition}
\begin{proof}
The two displayed substitutions compose to the identity in all three
coordinates, without a division. Substituting each $a_j^g$ into $t_0,t_1,t_2$
gives the three formulas. For example
$t_0(a_1^g)=-6g-y+6g+x+y-z+2g-x-2g=-z=t_2$;
$t_1(a_1^g)=-6g-y-z+2g-x+2g=t_0$ and
$t_2(a_1^g)=-z+2g-x=t_1$. This proves the first conjugacy;
the second and third follow by the same explicitly given substitutions.
The inverse coordinate formulas retain both $g$ and all integer coordinates.
\end{proof}

\begin{proposition}[The hexagonal quotient and its exact integral inverse]
Set $X=t_2-t_1$, $Y=t_0-t_1$, $W=-(t_0+t_1+t_2)$. The image in
$\mathbb Z^3$ is exactly
\[
 \mathcal L=\{(X,Y,W):W\equiv2X-Y\pmod3\},
\]
an index-three sublattice. Its inverse is
\[
 t_0={-W-X+2Y\over3},\quad t_1={-W-X-Y\over3},\quad
 t_2={-W+2X-Y\over3}.
\]
For $Q(X,Y)=X^2-XY+Y^2$, one has
$2Q(X,Y)+W^2=3(t_0^2+t_1^2+t_2^2)$.
The integral operation $r(t)=(-t_1,-t_2,-t_0)$ satisfies
\[
 r^2=c,\quad r^3=-I,\quad r^6=I,\qquad
 r:(X,Y,W)\longmapsto(X-Y,X,-W).
\]
It preserves $\mathcal L$ and $Q$. It is a specified extension of the cyclic
action, not claimed to be a word in the original monodromy group.
The second original monodromy acts on the planar readout by
$(X,Y)\mapsto(Y+2g,-X)$ and satisfies
\[
 Q(Y+2g,-X)-Q(X,Y)=2XY+2gX+4gY+4g^2.
\]
\end{proposition}
\begin{proof}
Solving the three linear equations gives the displayed inverse. Its three
numerators are divisible by three exactly under the one displayed congruence.
The map $(X,Y,W)\mapsto W-2X+Y\pmod3$ is onto, so its kernel has index three.
Expanding $2(X^2-XY+Y^2)+W^2$ after substitution cancels every cross term
and leaves $3\sum t_i^2$. Direct iteration gives the powers of $r$.
The new readouts are $X-Y,X,-W$; their gluing residual is
$-W-2(X-Y)+X=-W-X+2Y$, divisible by three on $\mathcal L$.
Expanding $Q(X-Y,X)$ gives $Q(X,Y)$. Substitution into $b_g$ gives
$X'=Y+2g,Y'=-X$; expanding the two quadratic polynomials proves the last
identity, including its nonzero correction.
\end{proof}
The readout $(t_0,t_1,t_2)\mapsto(X,Y)$ has kernel
$\mathbb Z(1,1,1)$. Retaining $W$ restores that omitted coordinate subject
to the exact gluing above. Modulo three the displayed divided inverse is
\emph{not} silently used: the original integral $t$-coordinates remain the
primary state at such levels. No unrestricted product of a plane and an
integer axis replaces $\mathcal L$.

\begin{proposition}[Gate order, compression and the simultaneous local return]
Put $d=\gcd(R,g)$ and $D=R/d$. For
$\Delta=\langle a_1,a_2:a_1^3=a_2^4=1\rangle$, write the above affine
representation as $v\mapsto B_jv+g b_j$. Then
\[
 \operatorname{ord}(g[b_R])=D
 \quad\hbox{in }H^1(\Delta,(\mathbb Z/R\mathbb Z)^3_B).
\]
Multiplication by $g$ gives an equivariant injection from the coefficient-one
system at level $D$ to the coefficient-$g$ system at level $R$:
\[
 (\mathbb Z/D\mathbb Z)^3\longrightarrow(\mathbb Z/R\mathbb Z)^3,
 \qquad v\longmapsto gv.
\]
Its image is $(d\mathbb Z/R\mathbb Z)^3$. Its inverse on the image divides
representatives by $d$ and multiplies by $(g/d)^{-1}$ modulo $D$.
Both systems have a global fixed sheet precisely when $D=1$. The full
level-$R$, coefficient-zero system then has $R$ fixed sheets, while the
compressed level-one system has one.
\end{proposition}
\begin{proof}
Block multiplication gives $b_{hk}=b_h+B_hb_k$. A coboundary at the cusp
has third coordinate zero because the third row of $B_\infty-I$ is zero.
Third cusp evaluation is consequently a well-defined homomorphism on
cohomology and sends $g[b_R]$ to $-g\pmod R$. If $k g[b_R]=0$, then
$R\mid kg$, equivalently $D\mid k$. Conversely $Dg=0\pmod R$ kills every
cocycle coefficient, proving exact order.
Since $\gcd(g/d,D)=1$, the stated inverse proves injectivity and identifies
the image. Each affine generator obeys
$g(B_jv+b_j)=B_j(gv)+g b_j$, proving equivariance. For $D=1$ every
division-modulo-one means its unique zero class.
A common fixed point would be fixed by the cusp; its last coordinate forces
$g=0\pmod R$. When $g=0$, the second generator forces $x=-y$ and $y=x$,
so $x=y=0$ since $R$ is odd. Every $z$ is then fixed by both generators.
This proves the fixed-sheet counts and their existence equivalence.
\end{proof}
In an ES leaf, $g=4u+1$ in channel E and $g=u+a$ in channel M. Thus the
actual subgroup $\langle[g]_R\rangle$ and the subgroup generated by
$g[b_R]$ are isomorphic by $kg\mapsto kg[b_R]$; minus cusp evaluation is
the inverse. For $R=\prod\ell^{h_\ell}$, the CRT map is componentwise
reduction of this \emph{one} labelled state. Its inverse is
$v=\sum_\ell v_\ell (R/\ell^{h_\ell})
 ((R/\ell^{h_\ell})^{-1}\bmod\ell^{h_\ell})\pmod R$ in each coordinate.
It retains the same divisor exponents, channel and gate at every factor.
The order is $D=\prod_\ell\ell^{\max(h_\ell-v_\ell(g),0)}$.
Different divisors chosen separately at different primes do not define this
same-leaf correspondence.

\section{JT's quartic and the positive-square locus of the hexagon}
Let $\eta=e^{\pi i/3}$ and $\zeta=e^{\pi i/6}$, so $\zeta^2=\eta$.
The field $\mathbb Q(\zeta)$ has degree four over $\mathbb Q$; its
chosen degree-two subfield is $\mathbb Q(\eta)$ and the nontrivial relative
automorphism sends $\zeta$ to $-\zeta$. For integers $r,s$,
\[
 N_{\mathbb Q(\zeta)/\mathbb Q(\eta)}(r-s\zeta)=r^2-\eta s^2,
 \qquad
 N_{\mathbb Q(\eta)/\mathbb Q}(X-\eta Y)=X^2-XY+Y^2.
\]
Indeed the first product is $(r-s\zeta)(r+s\zeta)$; the second uses
$\eta+\bar\eta=1$ and $\eta\bar\eta=1$. Consequently
\[
 r^4-r^2s^2+s^4=Q(r^2,s^2)
 =N_{\mathbb Q(\zeta)/\mathbb Q}(r-s\zeta).
\]
Multiplication by $\eta^{-1}=1-\eta$ sends
$X-\eta Y$ to $(X-Y)-\eta X$, the exact planar hexagon map above.

\begin{proposition}[All positive-square returns in the dihedral orbit]
Let $r,s$ be distinct positive coprime integers. Under the group generated
by $(X,Y)\mapsto(X-Y,X)$ and $(X,Y)\mapsto(Y,X)$, the points having
both coordinates positive squares correspond to $(r,s),(s,r)$, and,
exactly when $|r^2-s^2|=t^2$ for a positive integer $t$, the two additional
pairs $(\max(r,s),t),(t,\max(r,s))$. Their common quartic norm does not
change. These statements retain the group element as a mark whenever the
chosen path matters.
\end{proposition}
\begin{proof}
The six rotation images of $(X,Y)$ are
\[
 (X,Y),(X-Y,X),(-Y,X-Y),(-X,-Y),(Y-X,-X),(Y,Y-X).
\]
The other six images are their swaps. For distinct $X,Y>0$, the positive
pairs in this list are exactly
$(X,Y),(Y,X),(\max(X,Y),|X-Y|),(|X-Y|,\max(X,Y))$.
Taking $X=r^2,Y=s^2$ proves the square criterion and the complete list.
Any common prime divisor of $t$ and $\max(r,s)$ would divide the other
original coordinate too, contradicting $\gcd(r,s)=1$.
Rotation and swapping both preserve $Q$ by polynomial expansion.
\end{proof}
For $(r,s)=(3,5)$, $25-9=16$ gives $(3,5),(5,3),(4,5),(5,4)$,
all of norm $481=37\cdot13$. At $(1,26)$ the difference is $675$,
not a square since $25^2<675<26^2$. The norm is
$456301=2521\cdot181$, and $88^2=3\cdot2521+181$.
Thus the exact norm bridge preserves this cofactor, which is a quadratic
residue modulo 2521. A norm-changing arithmetic operation is not supplied
by the hexagonal unit action. This is the precise scope of the bridge,
not a conclusion that the constructions are unrelated.


\section{Local maps of the completed affine covers}
For this section $\mathcal F_D=(\mathbb Z/D\mathbb Z)^3$ with the
coefficient-one affine action displayed above. The base is the sphere
with its three marked punctures removed, with their positively oriented
monodromies $a_1,a_2,a_\infty$. A finite equivariant set defines its
cover by taking the universal cover times the set and identifying
$(\gamma z,v)$ with $(z,\gamma^{-1}v)$ under deck transformations.
The marked cycles at a puncture are filled by discs with base coordinate
$t=w^e$, one for each cycle of length $e$. This construction specifies
the objects and local coordinates used in the following proof.
For the cusp $a_\infty(x,y,z)=(x,y+x,z-1)$, its $k$th iterate is
$(x,y+kx,z-k)$; its cycle length at every point of level $D$ is exactly
$D$. At the finite points the lengths divide three and four by the
displayed generator relations. The original level-modulus quotient
alone therefore does not specify those finite local degrees.
% Copy-ready replacement for src/core.tex, Proposition Functorial modulus maps.
% This private correction does not alter the supplied archive.
\begin{proposition}[Functorial modulus maps and their compactifications]
\label{prop:crt}
For positive odd integers $D\mid E$, coordinate reduction
$\mathcal F_E\to\mathcal F_D$ is equivariant and surjective, with fibres
of size $(E/D)^3$. It induces a map of the corresponding compactified
covers. At a marked source branch cycle of length $e_E$ mapping to a
target branch cycle of length $e_D$, one has $e_D\mid e_E$, and the local
degree is $e_E/e_D$. At the cusp this is $E/D$.

For coprime positive odd integers $D,E$, coordinate CRT gives an
equivariant bijection
\[
 \mathcal F_{DE}\xrightarrow{\sim}\mathcal F_D\times\mathcal F_E,
\]
with coordinate inverse
$a+D((b-a)D^{-1}\bmod E)\bmod DE$.
The corresponding unramified cover is the fibre product over the
punctured base. Its compactification is the normalization of the fibre
product of the compactified covers over $\mathbb P^1$.
The cocycles and local systems have the induced reduction and CRT maps.
Every product and every map retains the same base and candidate label.
\end{proposition}
\begin{proof}
All three affine formulas have integral coefficients, hence commute with
coordinate reduction. Each residue modulo $D$ has exactly $E/D$ lifts
modulo $E$, proving the fibre count in three coordinates. Equivariance
gives the map of covers away from the punctures.

At a fixed puncture, let $e_E$ be the source cycle length and $e_D$ its
image cycle length. The $e_E$-th power of target monodromy fixes the
image sheet, so $e_D\mid e_E$. Choose compatible sheet origins and
local coordinates with base parameter
\[
 t=w_E^{e_E}=w_D^{e_D}.
\]
The prescribed punctured-cover map has the local form
\[
 w_D=\zeta w_E^{e_E/e_D},\qquad \zeta^{e_D}=1,
\]
where the root of unity retains the target branch mark. The positive
integer exponent proves holomorphic extension, unique by agreement on
the punctured disc. At infinity $e_E=E$ and $e_D=D$ by the cusp-cycle
theorem. At the finite branch points their values are the actual cycle
lengths dividing three and four; they need not equal the moduli.

The displayed CRT formula is inverse to both coordinate reductions.
It commutes with all three monodromies, so it identifies the unramified
level-$DE$ cover with the fibre product of the two unramified covers.
To complete this identification, consider source and target branch
cycles of lengths $e_D,e_E$. Their local fibre product has equation
\[
 x^{e_D}=y^{e_E}.
\]
Put $d=\gcd(e_D,e_E)$, $r=e_D/d$, $s=e_E/d$. Its $d$ branches are
$x^r=\xi y^s$ for $\xi^d=1$. On a branch, choose $\eta^r=\xi$.
The normalization has the parametrization
\[
 x=\eta w^s,\qquad y=w^r,\qquad
 t=w^{\operatorname{lcm}(e_D,e_E)}.
\]
The two maps to the base agree because $\eta^{e_D}=\xi^d=1$.
Since $\gcd(r,s)=1$, Bezout's identity gives a birational inverse in
the field of meromorphic germs, and the parameter disc is smooth.
The product monodromy has precisely $d$ cycles of length
$\operatorname{lcm}(e_D,e_E)$: a simultaneous iterate returns exactly
when divisible by both lengths, and the number of orbits is
$e_De_E/\operatorname{lcm}(e_D,e_E)=d$.
Thus the normalized local branches are exactly the completed product
cover. Their unique extensions glue to the compactified level-$DE$
cover. This argument is componentwise for disconnected covers and
keeps every candidate label and marked branch.

For example, reduction from level 15 to level 3 sends the $a_1$ cycle
\[
 (0,0,0)\mapsto(6,9,13)\mapsto(0,9,2)\mapsto(0,0,0)
\]
to the cycle
\[
 (0,0,0)\mapsto(0,0,1)\mapsto(0,0,2)\mapsto(0,0,0).
\]
The local degree here is $3/3=1$, whereas at the cusp it is $15/3=5$.
For the coprime levels three and five the cusp fibre product
$x^3=y^5$ is singular; its normalization is
$x=w^5,y=w^3,t=w^{15}$. This records why normalization is part of
the compactified product statement.
\end{proof}

\section{The marked cubic-surface contraction and the toric hexagon}
\label{sec:picard-hexagon}
This section reconstructs the explicit geometric comparison in the supplied
discussion, lines 2017--2055. All point labels and directions below are part
of the objects. The split degree-six model and its six boundary curves occur
in \href{https://www.math.uni-bielefeld.de/lag/man/290.pdf}{Mark Blunk,
\emph{Del Pezzo surfaces of degree 6 over an arbitrary field}, \S2,
Proposition 2.1, printed p.~3}; the Picard maps, signs and torus coordinates
needed for the present comparison are proved here.

\subsection{The marked surfaces, contraction and complete inverse fibres}
Work over $\mathbb C$. Let $p_1,\ldots,p_6$ be six distinct marked points
of $\mathbb P^2$, no three collinear and not all six on a conic. Choose
coordinates with
\[
 p_1=[1:0:0],\quad p_2=[0:1:0],\quad p_3=[0:0:1].
\]
This is a nonempty explicitly realized domain: one choice is
\[
 p_4=[1:1:1],\qquad p_5=[1:2:3],\qquad p_6=[1:3:2].             \tag{PH1}
\]
For this choice the determinants of the twenty point triples, in
lexicographic index order, are
\[
1,1,3,2,-1,-2,-3,1,-1,-5,1,1,1,-2,-1,1,1,2,1,-3,
\]
with the points as rows in increasing label order; in particular they are all
nonzero. A conic through the first three points has the form
$A x_0x_1+B x_0x_2+C x_1x_2=0$. Its values at the last three have coefficient
matrix
\[
\begin{pmatrix}1&1&1\\2&3&6\\3&2&6\end{pmatrix},\qquad\det=7,
\]
so no conic passes through all six.

Put
\[
Y=\operatorname{Bl}_{p_1,p_2,p_3}\mathbb P^2,
\qquad S=\operatorname{Bl}_{p_1,\ldots,p_6}\mathbb P^2.
\]
Let $q_a\in Y$ be the unchanged point above $p_a$, for $a=4,5,6$.
All three lie in the complement of the coordinate triangle and its
exceptional divisors. The morphism
\[
 f:S=\operatorname{Bl}_{q_4,q_5,q_6}Y\longrightarrow Y             \tag{PH2}
\]
contracts exactly the three exceptional curves $E_4,E_5,E_6$.
Denote the line pullback and exceptional classes on $S$ by $H,E_1,\ldots,E_6$,
and those on $Y$ by $h,e_1,e_2,e_3$.

Here are the fibres and their inverse data explicitly. Near a centre $q_a$
take regular local coordinates $(\xi,\eta)$ vanishing there. The blowup is
\[
 \{(\xi,\eta,[s:t]):\xi t=\eta s\}.
\]
Its two charts are $(\xi,t)$ with $\eta=\xi t$, and $(s,\eta)$ with
$\xi=s\eta$. Over the centre the fibre is the whole $\mathbb P^1$ of
directions $[s:t]$; elsewhere the unique direction is $[\xi:\eta]$.
Thus $f$ is an isomorphism away from the three specified fibres, and
retaining the centres and these exceptional directions gives the full
blowup return. No inverse at a contracted point discards its direction.

The Picard groups and intersection forms are
\[
\operatorname{Pic}(S)=\mathbb ZH\oplus\bigoplus_{i=1}^6\mathbb ZE_i,
\quad H^2=1,\quad H E_i=0,\quad E_iE_j=-\delta_{ij},
\]
and the analogous basis $(h,e_1,e_2,e_3)$ on $Y$. To see the additional
summand at a blowup, push a divisor to the smooth base and subtract its
total transform; the difference is supported on the exceptional curve.
Its coefficient is unique since that curve has self-intersection $-1$,
as follows from the transition of its normal coordinate between the two
charts above. Applying this argument successively gives the displayed bases.
The maps are exactly
\[
 f^*h=H,\quad f^*e_i=E_i\ (1\le i\le3),\qquad
 f_*H=h,\quad f_*E_i=e_i\ (i\le3),\quad f_*E_a=0\ (a\ge4).
                                                               \tag{PH3}
\]
Hence $f_*f^*=1$, the kernel of $f_*$ is
$\mathbb ZE_4\oplus\mathbb ZE_5\oplus\mathbb ZE_6$, and it is orthogonal
to $f^*\operatorname{Pic}(Y)$. For
$D=dH-\sum m_iE_i$, $D'=d'H-\sum m'_iE_i$, one has
\[
 f^*f_*D=D+\sum_{a=4}^6m_aE_a,\qquad
 (f_*D)(f_*D')=DD'+\sum_{a=4}^6m_am'_a.                         \tag{PH4}
\]
These identities prove both the precise lost sublattice and the correction
to intersection numbers. The canonical classes obey
\[
 -K_S=3H-\sum_{i=1}^6E_i=f^*(-K_Y)-E_4-E_5-E_6,
 \quad -K_Y=3h-e_1-e_2-e_3,
 \quad K_S^2=3,\quad K_Y^2=6.                                  \tag{PH5}
\]
The canonical-class formula follows in a blowup chart from
$d\xi\wedge d\eta=\xi\,d\xi\wedge dt$ and then from the canonical
class $-3h$ of the plane.

For completeness, these marked blowups give the stated cubic model.
The cubic forms through the six points have dimension four: their six
evaluation conditions are independent, since a conic through any five
points, multiplied by a line avoiding the sixth, separates that sixth
evaluation. Such a conic exists, is unique by Bezout, and is irreducible
because a pair of lines contains at most four of the points. The cubics
have no further base point: for a point outside the six, the pairs whose
joining line contains it form a matching of at most three edges. There are
fifteen perfect matchings of six labels, and each forbidden edge belongs
to three; hence at least six perfect matchings avoid all forbidden edges.
Such a matching supplies a product of three joining lines
nonzero there. At each marked point, matchings with two different partners
give independent first-order directions. Thus the induced anticanonical
system is base-point free on $S$.
It is ample as well. For a nonexceptional irreducible curve of degree $d$
and multiplicities $m_i$, distinct from the six conics above, Bezout with
each such conic gives $2d\ge\sum_{j\ne i}m_j$. Summing gives
$\sum m_i\le12d/5<3d$, so $(-K_S)C>0$. Each of the six conics and every
$E_i$ has anticanonical degree one. Together with $K_S^2=3$, the
Nakai--Moishezon criterion proves ampleness.
The resulting morphism to $\mathbb P^3$ is finite, and its nondegenerate
surface image has degree at least two. The equality
$3=\deg(\text{morphism})\deg(\text{image})$ makes it birational onto a
cubic surface. A general hyperplane pullback is a smooth curve of genus
one by Bertini and adjunction. Its plane cubic image has arithmetic genus
one, so it is smooth; hence the cubic surface has no singular curve.
A hypersurface with no codimension-one singularity is normal. The finite
birational morphism to this normal surface is an isomorphism. In particular
the anticanonical model is a smooth cubic surface.

\subsection{All twenty-seven line classes and their contracted images}
The complete list on $S$ is
\[
 E_i\ (1\le i\le6),\qquad L_{ij}=H-E_i-E_j\ (i<j),\qquad
 Q_i=2H-\sum_{j\ne i}E_j\ (1\le i\le6).                        \tag{PH6}
\]
The last two classes are the strict transforms of the actual joining lines
and the unique conics through the indicated five points. Their curves are
smooth rational, have square $-1$, and anticanonical degree one, so their
images in the cubic model are lines. Conversely a line is a smooth rational
curve with square $-1$ by adjunction. Unless exceptional, it has degree
$d\ge1$ and nonnegative multiplicities satisfying
\[
 \sum m_i=3d-1,\qquad \sum m_i^2=d^2+1.
\]
Cauchy gives $(3d-1)^2\le6(d^2+1)$, hence $d\le2$. At $d=1$ exactly two
multiplicities are one; at $d=2$ exactly five are one. This proves (PH6)
is exhaustive without omitting a line class.

Let $J=\{1,2,3\}$ and $K=\{4,5,6\}$. The entire contraction correspondence
is the following table. A zero divisor class in its second row records a
contracted curve, not an inverse point of a divisor map.
\[
\begin{array}{c|c|c}
\text{curve on }S&\text{image class on }Y&\text{image curve square}\\\hline
E_i\ (i\in J)&e_i&-1\\
E_a\ (a\in K)&0\quad(\text{point }q_a)&\text{not a curve}\\
L_{ij}\ (i,j\in J)&h-e_i-e_j&-1\\
L_{ia}\ (i\in J,a\in K)&h-e_i&0\\
L_{ab}\ (a,b\in K)&h&1\\
Q_i\ (i\in J)&2h-\sum_{j\in J\setminus\{i\}}e_j&2\\
Q_a\ (a\in K)&2h-e_1-e_2-e_3&1
\end{array}                                                     \tag{PH7}
\]
Each noncontracted entry is the actual line or conic with its remaining
marked incidence conditions; coincident divisor classes do not identify
distinct image curves. The total-transform inverse is
\[
 f^*f_*L_{ij}=L_{ij}+\sum_{a\in K\cap\{i,j\}}E_a,
 \qquad f^*f_*Q_i=Q_i+\sum_{a\in K\setminus\{i\}}E_a,           \tag{PH8}
\]
because the original plane curves pass simply through exactly those centres.
Equations (PH3)--(PH4) prove the table and every square in it.

The six curves
\[
 E_1,L_{12},E_2,L_{23},E_3,L_{13}                                \tag{PH9}
\]
avoid the three contraction centres and return isomorphically to their six
namesakes on $Y$. Each has square $-1$. The intersection rules give
$E_iL_{jk}=1$ exactly when $i\in\{j,k\}$; different $E_i$ are disjoint,
and the three displayed $L_{ij}$ are pairwise disjoint. Thus (PH9) is
exactly the cycle with neighbouring intersection one and every other
intersection zero. Its sum is $-K_Y$.

\subsection{A six-cycle realized by Cremona and a marked configuration map}
Put $\alpha=H-E_1-E_2-E_3$, so $\alpha^2=-2$ and $K_S\alpha=0$.
Define the reflection and permutation
\[
 C(D)=D+(D\alpha)\alpha,
 \qquad \sigma=(1\ 3\ 2),\quad
 \sigma(H)=H,\quad\sigma(E_i)=E_{\sigma(i)},
\]
where $\sigma$ fixes the labels in $K$. Their composite is $T=\sigma C$.
The two operations commute, $C^2=1$, and $\sigma^3=1$. Expansion using
$\alpha^2=-2$ proves $C(D)C(D')=DD'$, and both operations fix $K_S$.
Consequently $T$ is an integral intersection isometry, fixes the canonical
class, and $T^3=C$, $T^6=1$. Its exact basis action is
\[
\begin{split}
T(H)&=2H-E_1-E_2-E_3,\\
T(E_1)&=L_{12},\quad T(E_2)=L_{23},\quad T(E_3)=L_{13},\qquad
T(E_a)=E_a\ (a\in K).
\end{split}                                                     \tag{PH10}
\]
In particular it sends the six entries of (PH9) successively to the next,
including $L_{13}\mapsto E_1$, and its order is exactly six.
Its full permutation of all twenty-seven classes is
\[
\begin{gathered}
(E_1\ L_{12}\ E_2\ L_{23}\ E_3\ L_{13}),\qquad
(E_4)(E_5)(E_6),\\
(L_{14}\ L_{34}\ L_{24})\ (L_{15}\ L_{35}\ L_{25})\
(L_{16}\ L_{36}\ L_{26}),\qquad (Q_1\ Q_3\ Q_2),\\
(L_{45}\ Q_6)\ (L_{46}\ Q_5)\ (L_{56}\ Q_4).
\end{gathered}                                                   \tag{PH11}
\]
To verify completeness directly, a mixed class $L_{ia}$ has $L_{ia}\alpha=0$,
so maps to $L_{\sigma(i),a}$; a class $L_{ab}$ in $K$ has product one with
$\alpha$, so maps to $Q_c$ for $K\setminus\{a,b\}=\{c\}$; the latter
has product $-1$ and maps back. A $Q_i$ with $i\in J$ has product zero
and maps to $Q_{\sigma(i)}$. Together with (PH10), these calculations prove
every entry of (PH11); intersection preservation holds for every pair,
not only for neighbouring sides of the hexagon.

This is realized geometrically on the degree-six surface. The quadratic
plane Cremona map and the coordinate permutation are
\[
 \kappa[x_0:x_1:x_2]=[x_1x_2:x_0x_2:x_0x_1],\qquad
 \rho[x_0:x_1:x_2]=[x_1:x_2:x_0].
\]
Their composite is
\[
 F=\rho\kappa:[x_0:x_1:x_2]\dashrightarrow
 [x_0x_2:x_0x_1:x_1x_2].                                       \tag{PH12}
\]
The blowup $Y$ is also the graph closure
\[
 \{([x],[y])\in\mathbb P^2\times\mathbb P^2:
                 x_0y_0=x_1y_1=x_2y_2\}.
\]
Away from the three coordinate points the second coordinate is the
displayed Cremona value; at each point the two blowup charts give its
exceptional direction. Thus switching the two projective factors extends
$\kappa$ to an automorphism of $Y$, and simultaneous permutation of the
coordinates extends $\rho$. Their action on curve classes is the
\emph{pushforward} $T=\sigma C$; the pullback is $T^{-1}=C\sigma^{-1}$.
For instance a general line has Cremona transform a conic through the three
centres, and $E_i$ maps to the opposite joining line, proving (PH10) directly.

On a six-point blowup the exact geometric statement retains the changed
centres. Set $q'_a=F(q_a)$ and
$S'=\operatorname{Bl}_{q'_4,q'_5,q'_6}Y$. Then $F$ lifts to an isomorphism
$\widetilde F:S\to S'$ making $f'\widetilde F=Ff$, with the exceptional
label $a$ sent to the exceptional label $a$ at $q'_a$. Its inverse is the
lift of $F^{-1}$. On exceptional directions its map is the projectivized
invertible derivative $\mathbb P(dF_{q_a})$, obtained by substituting the
two local blowup charts. For (PH1), the new plane centres are exactly
\[
 p'_4=[1:1:1],\qquad p'_5=[3:2:6],\qquad p'_6=[2:3:6].
\]
Thus (PH10)--(PH11) are the complete marked pushforward between these
configurations; they do not assert an automorphism of the fixed generic
six-point configuration. The contraction maps commute with $T$ on their
Picard groups because the kernel spanned by $E_4,E_5,E_6$ is fixed.

\subsection{Torus coordinates, the exact fan, and the cusp-side convention}
Use the following specified coordinates on the dense torus of $Y$:
\[
 u=x_2/x_0,\qquad v=x_2/x_1,
 \qquad [x_0:x_1:x_2]=[v:u:uv].                                \tag{PH13}
\]
Substitution in (PH12) gives the genuine torus automorphism
\[
 F(u,v)=(u/v,u),\qquad F^{-1}(u,v)=(v,v/u),\qquad
 F^3(u,v)=(u^{-1},v^{-1}),\quad F^6(u,v)=(u,v).                  \tag{PH14}
\]
For example $F^2(u,v)=(v^{-1},u/v)$, proving the third-power formula by
one more substitution. At a retained centre $(u,v)$ the exceptional
direction map in (PH2) is completely explicit:
\[
dF_{(u,v)}=\begin{pmatrix}1/v&-u/v^2\\1&0\end{pmatrix},\quad
\det dF=u/v^2,\qquad
[P:Q]\longmapsto[vP-uQ:v^2P].                                  \tag{PH14a}
\]
At the image centre $(U,V)=(u/v,u)$ its inverse sends
$[P':Q']$ to $[U^2Q':UQ'-VP']$. Substitution gives
$[u^2P:u^2Q]=[P:Q]$, proving the inverse on the entire exceptional
$\mathbb P^1$, including vertical and infinite slope directions.
The valuation vectors
$(\operatorname{ord}_B u,\operatorname{ord}_B v)$ of the boundary curves
in the retained order are precisely
\[
\begin{array}{c|rrrrrr}
B&E_1&L_{12}&E_2&L_{23}&E_3&L_{13}\\\hline
\operatorname{ord}_B u&1&1&0&-1&-1&0\\
\operatorname{ord}_B v&0&1&1&0&-1&-1
\end{array}                                                     \tag{PH15}
\]
For a joining line these are its zero and pole orders in (PH13). At $E_1$,
both $x_1$ and $x_2$ vanish to first order while $x_0$ is a unit, giving
$(1,0)$; the other two exceptional rows follow in the same way from the
displayed coordinate ratios. Thus the six rays, in the exact curve order,
are
\[
n_1=(1,0),\ n_2=(1,1),\ n_3=(0,1),\ n_4=(-1,0),\
n_5=(-1,-1),\ n_6=(0,-1).
\]
Adjacent pairs have determinant one. The cones between them are the six
smooth toric charts of the three-point blowup: the original plane fan's
three cones are each subdivided by the sum of their two generating rays,
which is the blowup operation in the local charts of (PH2).
On this cocharacter lattice (PH14) acts by
\[
 R_6=\begin{pmatrix}1&-1\\1&0\end{pmatrix},\qquad
 R_6n_i=n_{i+1},\quad R_6^3=-I,
\]
with subscripts modulo six. Its action on characters is the transpose.
The equality $Q(X-Y,X)=Q(X,Y)$ for $Q=X^2-XY+Y^2$ proves that these are
the same six norm-one vectors as in the integral hexagonal quotient.
Because $R_6$ carries every cone to the next cone, its monomial map and
inverse carry the dual-cone monomial coordinate rings into the corresponding
chart rings. This proves extension across every boundary chart directly,
in agreement with the graph model above.

The source cusp marks its rays, in \S9.3 of the
\href{https://alpo.ge/s6.pdf}{circulated period manuscript}, by
\[
(1,0),(1,-1),(0,-1),(-1,0),(-1,1),(0,1).
\]
The integral matrix $B=\operatorname{diag}(1,-1)$ sends each $n_i$ to the
corresponding source ray, preserves opposite pairs, and is its own inverse.
It induces the exact torus coordinate change $(u,v)\mapsto(u,v^{-1})$.
In that convention the conjugate rotation is
$BR_6B=\left(\begin{smallmatrix}1&1\\-1&0\end{smallmatrix}\right)$ and
the conjugate torus map is $(u,v)\mapsto(uv,u^{-1})$.

This identifies the marked smooth toric normalizations and their six
curves. The source further supplies a normalization map $\nu:dP_6\to W$
identifying each side with its opposite. With that source map retained,
the actual comparison is the composite $S\xrightarrow{f}Y\xrightarrow{B}
dP_6\xrightarrow{\nu}W$. The first map has exactly the three
$\mathbb P^1$ fibres proved above. The last has one preimage off the
boundary, two branches at a general double-curve point, and three at each
triple point, as its local equations $xy=0$ and $xyz=0$ show. The source
marks the two triple-point preimages by the alternating vertex sets
$\{p_{12},p_{34},p_{56}\}$ and $\{p_{23},p_{45},p_{61}\}$.
All side-identification maps remain attached to $\nu$; the fan map alone
does not reconstruct their gluing parameters or identify the singular
quotient with its smooth normalization. No global six-sphere claim is used
in the Picard, torus or contraction proofs.


\begin{thebibliography}{9}
\bibitem{blunk} Mark Blunk, \emph{Del Pezzo surfaces of degree 6 over an
arbitrary field}, arXiv:0805.0119v2, Section 2, Proposition 2.1.
\url{https://arxiv.org/abs/0805.0119v2}.
\bibitem{jt} JT and The Clankers, \emph{Boolean product formulation and
joint ES research discussion}, September 2026. Complete supplied research
dialogues, exact export identities and statement locators in
\texttt{../PROVENANCE.md}. These are research provenance, not a substitute
for the proofs written here.
\bibitem{et} Christian Elsholtz and Terence Tao, \emph{Counting the number
of solutions to the Erd\H{o}s--Straus equation on unit fractions},
J. Aust. Math. Soc. \textbf{94} (2013), 50--105, Section 2.
\url{https://arxiv.org/abs/1107.1010}.
\bibitem{period} \emph{The $(3,4,\infty)$ modular family of 2-tori,
completed at its three special points, is a complex structure on $S^6$},
108-page circulated manuscript hosted by Levent Alp\"oge,
\url{https://alpo.ge/s6.pdf}. Citation identifies the source of the period
matrices, not a certification of the title's conclusion or an inference
of authorship from hosting. See the received bibliography for the
accompanying Engel exposition and attribution discussion.
\bibitem{sv} Jeroen Sijsling and John Voight, \emph{On computing Bely\u\i
maps}, Publications math\'ematiques de Besan\c con 2014/1, 73--131,
Section 1. \url{https://arxiv.org/abs/1311.2529}.
\end{thebibliography}
\end{document}
